%%%%%%%%%%%%%%%%%%%%%%%%%%%%%%%%%%%%%%%%%%%%%%%%%%%%%%%%%%%%
% 13.8 MATHEMATICAL FINANCE
% The Composition of Advanced Mathematics
%%%%%%%%%%%%%%%%%%%%%%%%%%%%%%%%%%%%%%%%%%%%%%%%%%%%%%%%%%%%

\section{Mathematical Finance}
\label{sec:mathematical-finance}

\prerequisite{
Probability, statistics, stochastic processes, stochastic calculus,
differential equations, optimization, linear algebra, numerical methods,
and mathematical modeling.
}

\learningobjectives{
By the end of this section, the reader should be able to:
\begin{itemize}
    \item explain the purpose of mathematical finance;
    \item distinguish assets, securities, portfolios, and derivatives;
    \item compute simple and continuously compounded returns;
    \item distinguish arithmetic and logarithmic returns;
    \item understand discounting and compounding;
    \item compute present and future values;
    \item understand discount factors;
    \item recognize yield and interest-rate conventions;
    \item understand bonds and zero-coupon securities;
    \item compute bond prices from discounted cash flows;
    \item understand duration and convexity conceptually;
    \item define financial risk and return;
    \item compute portfolio expected return;
    \item compute portfolio variance and covariance;
    \item understand diversification;
    \item formulate mean-variance portfolio optimization;
    \item recognize the efficient frontier;
    \item understand minimum-variance portfolios;
    \item understand the role of correlations in diversification;
    \item recognize short-sale constraints;
    \item understand risk-free assets;
    \item recognize Sharpe ratios;
    \item understand market and idiosyncratic risk conceptually;
    \item recognize factor models;
    \item understand the Capital Asset Pricing Model conceptually;
    \item define beta;
    \item understand arbitrage;
    \item recognize the law of one price;
    \item understand replication;
    \item recognize complete and incomplete markets;
    \item understand contingent claims;
    \item define European call and put options;
    \item understand option payoffs;
    \item derive put-call parity;
    \item recognize forward and futures contracts;
    \item compute forward prices under simple assumptions;
    \item understand geometric Brownian motion;
    \item recognize stochastic differential equations for asset prices;
    \item apply It\^o's formula conceptually;
    \item understand risk-neutral valuation;
    \item recognize equivalent martingale measures conceptually;
    \item understand discounted asset prices under risk-neutral measures;
    \item derive the Black--Scholes equation conceptually;
    \item recognize the Black--Scholes option formula;
    \item understand option Greeks;
    \item interpret delta, gamma, theta, vega, and rho;
    \item understand delta hedging;
    \item recognize volatility;
    \item distinguish historical and implied volatility;
    \item understand volatility smiles and surfaces conceptually;
    \item recognize binomial option-pricing models;
    \item understand backward induction;
    \item price basic derivatives using binomial trees;
    \item understand American options conceptually;
    \item recognize optimal early exercise;
    \item understand Monte Carlo pricing;
    \item recognize variance reduction conceptually;
    \item understand finite-difference pricing conceptually;
    \item recognize stochastic volatility models conceptually;
    \item understand jump-diffusion models conceptually;
    \item recognize interest-rate models conceptually;
    \item understand term structures and yield curves;
    \item distinguish spot and forward rates;
    \item understand credit risk conceptually;
    \item recognize default probabilities and recovery rates;
    \item understand Value at Risk;
    \item recognize limitations of Value at Risk;
    \item understand Conditional Value at Risk;
    \item formulate financial optimization under risk constraints;
    \item understand dynamic portfolio choice conceptually;
    \item recognize stochastic control in finance;
    \item understand transaction costs;
    \item recognize liquidity risk;
    \item understand market impact conceptually;
    \item recognize high-frequency data issues conceptually;
    \item understand calibration and parameter estimation;
    \item recognize model risk;
    \item distinguish statistical estimation error from financial model error;
    \item understand scenario and stress testing;
    \item recognize systemic financial networks conceptually;
    \item connect mathematical finance with probability, statistics,
          optimization, stochastic calculus, PDEs, numerical analysis,
          control theory, and network science.
\end{itemize}
}

%%%%%%%%%%%%%%%%%%%%%%%%%%%%%%%%%%%%%%%%%%%%%%%%%%%%%%%%%%%%
\subsection{What Is Mathematical Finance?}
\label{subsec:what-is-mathematical-finance}
%%%%%%%%%%%%%%%%%%%%%%%%%%%%%%%%%%%%%%%%%%%%%%%%%%%%%%%%%%%%

Mathematical finance uses mathematical models to study prices, risk,
investment decisions, and financial contracts.

Important mathematical tools include

\[
\boxed{
\text{probability}
}
\]

\[
\boxed{
\text{statistics}
}
\]

\[
\boxed{
\text{stochastic processes}
}
\]

\[
\boxed{
\text{optimization}
}
\]

\[
\boxed{
\text{differential equations}
}
\]

and

\[
\boxed{
\text{numerical computation}.
}
\]

The subject combines descriptive models of financial markets with
decision-making models for investment and risk management.

%%%%%%%%%%%%%%%%%%%%%%%%%%%%%%%%%%%%%%%%%%%%%%%%%%%%%%%%%%%%
\subsection{Financial Assets}
\label{subsec:financial-assets}
%%%%%%%%%%%%%%%%%%%%%%%%%%%%%%%%%%%%%%%%%%%%%%%%%%%%%%%%%%%%

A financial asset is a claim with economic value.

Examples include:

\[
\boxed{
\text{stocks},
}
\]

\[
\boxed{
\text{bonds},
}
\]

\[
\boxed{
\text{currencies},
}
\]

\[
\boxed{
\text{commodities},
}
\]

and

\[
\boxed{
\text{derivatives}.
}
\]

An asset price at time \(t\) will often be denoted

\[
\boxed{
S_t.
}
\]

%%%%%%%%%%%%%%%%%%%%%%%%%%%%%%%%%%%%%%%%%%%%%%%%%%%%%%%%%%%%
\subsection{Returns}
\label{subsec:financial-returns}
%%%%%%%%%%%%%%%%%%%%%%%%%%%%%%%%%%%%%%%%%%%%%%%%%%%%%%%%%%%%

Suppose an asset price changes from \(S_t\) to \(S_{t+1}\).

The simple return is

\[
\boxed{
R_{t+1}
=
\frac{S_{t+1}-S_t}{S_t}.
}
\]

Equivalently,

\[
\boxed{
1+R_{t+1}
=
\frac{S_{t+1}}{S_t}.
}
\]

%%%%%%%%%%%%%%%%%%%%%%%%%%%%%%%%%%%%%%%%%%%%%%%%%%%%%%%%%%%%
\subsection{Log Returns}
\label{subsec:log-returns}
%%%%%%%%%%%%%%%%%%%%%%%%%%%%%%%%%%%%%%%%%%%%%%%%%%%%%%%%%%%%

The continuously compounded or logarithmic return is

\[
\boxed{
r_{t+1}
=
\ln
\left(
\frac{S_{t+1}}{S_t}
\right).
}
\]

Thus

\[
\boxed{
r_{t+1}
=
\ln(1+R_{t+1}).
}
\]

Log returns are additive across time:

\[
\boxed{
\ln\frac{S_T}{S_0}
=
\sum_{t=0}^{T-1}
\ln\frac{S_{t+1}}{S_t}.
}
\]

%%%%%%%%%%%%%%%%%%%%%%%%%%%%%%%%%%%%%%%%%%%%%%%%%%%%%%%%%%%%
\subsection{Worked Example: Returns}
\label{subsec:worked-financial-returns}
%%%%%%%%%%%%%%%%%%%%%%%%%%%%%%%%%%%%%%%%%%%%%%%%%%%%%%%%%%%%

\begin{workedexamplebox}{Simple and Log Return}

Suppose a stock increases from

\[
S_0=100
\]

to

\[
S_1=108.
\]

The simple return is

\[
R
=
\frac{108-100}{100}
=
0.08.
\]

The logarithmic return is

\[
r
=
\ln\left(\frac{108}{100}\right).
\]

Therefore,

\begin{answerbox}
\[
\boxed{
R=8\%,
\qquad
r\approx0.07696.
}
\]
\end{answerbox}

\end{workedexamplebox}

%%%%%%%%%%%%%%%%%%%%%%%%%%%%%%%%%%%%%%%%%%%%%%%%%%%%%%%%%%%%
\subsection{Compounding}
\label{subsec:financial-compounding}
%%%%%%%%%%%%%%%%%%%%%%%%%%%%%%%%%%%%%%%%%%%%%%%%%%%%%%%%%%%%

Suppose an initial amount \(P\) earns interest rate \(r\) per period.

After \(n\) periods,

\[
\boxed{
A_n
=
P(1+r)^n.
}
\]

This is discrete compounding.

%%%%%%%%%%%%%%%%%%%%%%%%%%%%%%%%%%%%%%%%%%%%%%%%%%%%%%%%%%%%
\subsection{Continuous Compounding}
\label{subsec:continuous-compounding}
%%%%%%%%%%%%%%%%%%%%%%%%%%%%%%%%%%%%%%%%%%%%%%%%%%%%%%%%%%%%

For continuously compounded rate \(r\),

\[
\boxed{
A(t)
=
Pe^{rt}.
}
\]

The continuously compounded discount factor over time \(T\) is

\[
\boxed{
e^{-rT}.
}
\]

%%%%%%%%%%%%%%%%%%%%%%%%%%%%%%%%%%%%%%%%%%%%%%%%%%%%%%%%%%%%
\subsection{Present Value}
\label{subsec:present-value}
%%%%%%%%%%%%%%%%%%%%%%%%%%%%%%%%%%%%%%%%%%%%%%%%%%%%%%%%%%%%

If a payment \(C\) will be received at time \(T\), its present value under
constant continuously compounded rate \(r\) is

\[
\boxed{
PV
=
Ce^{-rT}.
}
\]

Present value converts future cash flows to current monetary units.

%%%%%%%%%%%%%%%%%%%%%%%%%%%%%%%%%%%%%%%%%%%%%%%%%%%%%%%%%%%%
\subsection{Future Value}
\label{subsec:future-value}
%%%%%%%%%%%%%%%%%%%%%%%%%%%%%%%%%%%%%%%%%%%%%%%%%%%%%%%%%%%%

A present amount \(P\) has future value

\[
\boxed{
FV
=
Pe^{rT}
}
\]

under continuous compounding.

Thus discounting and compounding are inverse operations.

%%%%%%%%%%%%%%%%%%%%%%%%%%%%%%%%%%%%%%%%%%%%%%%%%%%%%%%%%%%%
\subsection{Worked Example: Present Value}
\label{subsec:worked-present-value}
%%%%%%%%%%%%%%%%%%%%%%%%%%%%%%%%%%%%%%%%%%%%%%%%%%%%%%%%%%%%

\begin{workedexamplebox}{Discounting a Future Payment}

Suppose

\[
C=1000,
\qquad
r=0.05,
\qquad
T=3.
\]

Then

\[
PV
=
1000e^{-0.05(3)}.
\]

Hence

\[
PV
\approx
860.71.
\]

\begin{answerbox}
\[
\boxed{
PV\approx \$860.71.
}
\]
\end{answerbox}

\end{workedexamplebox}

%%%%%%%%%%%%%%%%%%%%%%%%%%%%%%%%%%%%%%%%%%%%%%%%%%%%%%%%%%%%
\subsection{Discount Factors}
\label{subsec:discount-factors}
%%%%%%%%%%%%%%%%%%%%%%%%%%%%%%%%%%%%%%%%%%%%%%%%%%%%%%%%%%%%

A discount factor from time \(T\) to time \(0\) may be written

\[
\boxed{
D(0,T).
}
\]

Under a constant continuously compounded rate,

\[
\boxed{
D(0,T)
=
e^{-rT}.
}
\]

With a term structure of rates, the discount factor can depend nontrivially on
maturity.

%%%%%%%%%%%%%%%%%%%%%%%%%%%%%%%%%%%%%%%%%%%%%%%%%%%%%%%%%%%%
\subsection{Zero-Coupon Bonds}
\label{subsec:zero-coupon-bonds}
%%%%%%%%%%%%%%%%%%%%%%%%%%%%%%%%%%%%%%%%%%%%%%%%%%%%%%%%%%%%

A zero-coupon bond pays a single amount at maturity.

If the face value is \(F\), then

\[
\boxed{
P(0,T)
=
F D(0,T).
}
\]

Under constant rate \(r\),

\[
\boxed{
P(0,T)
=
Fe^{-rT}.
}
\]

%%%%%%%%%%%%%%%%%%%%%%%%%%%%%%%%%%%%%%%%%%%%%%%%%%%%%%%%%%%%
\subsection{Coupon Bonds}
\label{subsec:coupon-bonds}
%%%%%%%%%%%%%%%%%%%%%%%%%%%%%%%%%%%%%%%%%%%%%%%%%%%%%%%%%%%%

A coupon bond pays a sequence of cash flows.

If payments are \(C_1,\ldots,C_n\), then

\[
\boxed{
P
=
\sum_{k=1}^{n}
C_kD(0,t_k).
}
\]

The final payment may include both coupon and principal.

Bond pricing is therefore a discounted-cash-flow problem.

%%%%%%%%%%%%%%%%%%%%%%%%%%%%%%%%%%%%%%%%%%%%%%%%%%%%%%%%%%%%
\subsection{Yield to Maturity}
\label{subsec:yield-to-maturity}
%%%%%%%%%%%%%%%%%%%%%%%%%%%%%%%%%%%%%%%%%%%%%%%%%%%%%%%%%%%%

Yield to maturity is the single discount rate \(y\) that satisfies

\[
\boxed{
P
=
\sum_{k=1}^{n}
\frac{C_k}{(1+y)^{t_k}}
}
\]

under the specified compounding convention.

It summarizes the bond's price as an implied rate.

For multiple cash flows, yield generally must be solved numerically.

%%%%%%%%%%%%%%%%%%%%%%%%%%%%%%%%%%%%%%%%%%%%%%%%%%%%%%%%%%%%
\subsection{Interest-Rate Sensitivity}
\label{subsec:interest-rate-sensitivity}
%%%%%%%%%%%%%%%%%%%%%%%%%%%%%%%%%%%%%%%%%%%%%%%%%%%%%%%%%%%%

Bond prices and yields move inversely.

When yields increase, discounted future cash flows become less valuable.

Thus

\[
\boxed{
\text{yield rises}
\Rightarrow
\text{bond price falls}
}
\]

under ordinary fixed-cash-flow conditions.

%%%%%%%%%%%%%%%%%%%%%%%%%%%%%%%%%%%%%%%%%%%%%%%%%%%%%%%%%%%%
\subsection{Duration}
\label{subsec:bond-duration}
%%%%%%%%%%%%%%%%%%%%%%%%%%%%%%%%%%%%%%%%%%%%%%%%%%%%%%%%%%%%

Duration measures sensitivity of bond price to changes in yield.

Modified duration approximately satisfies

\[
\boxed{
\frac{\Delta P}{P}
\approx
-D_{\mathrm{mod}}\Delta y.
}
\]

Duration is therefore a first-order interest-rate sensitivity measure.

%%%%%%%%%%%%%%%%%%%%%%%%%%%%%%%%%%%%%%%%%%%%%%%%%%%%%%%%%%%%
\subsection{Convexity}
\label{subsec:bond-convexity}
%%%%%%%%%%%%%%%%%%%%%%%%%%%%%%%%%%%%%%%%%%%%%%%%%%%%%%%%%%%%

Convexity captures curvature in the price-yield relationship.

A second-order approximation has form

\[
\boxed{
\frac{\Delta P}{P}
\approx
-D_{\mathrm{mod}}\Delta y
+
\frac12C(\Delta y)^2.
}
\]

Duration gives the local slope.

Convexity improves the approximation for larger yield changes.

%%%%%%%%%%%%%%%%%%%%%%%%%%%%%%%%%%%%%%%%%%%%%%%%%%%%%%%%%%%%
\subsection{Portfolios}
\label{subsec:financial-portfolios}
%%%%%%%%%%%%%%%%%%%%%%%%%%%%%%%%%%%%%%%%%%%%%%%%%%%%%%%%%%%%

A portfolio combines multiple assets.

Let

\[
w_i
\]

be the fraction invested in asset \(i\).

For a fully invested portfolio,

\[
\boxed{
\sum_{i=1}^{n}w_i=1.
}
\]

If short selling is prohibited,

\[
\boxed{
w_i\geq0.
}
\]

%%%%%%%%%%%%%%%%%%%%%%%%%%%%%%%%%%%%%%%%%%%%%%%%%%%%%%%%%%%%
\subsection{Portfolio Return}
\label{subsec:portfolio-return}
%%%%%%%%%%%%%%%%%%%%%%%%%%%%%%%%%%%%%%%%%%%%%%%%%%%%%%%%%%%%

If asset returns are \(R_1,\ldots,R_n\), then portfolio return is

\[
\boxed{
R_p
=
\sum_{i=1}^{n}
w_iR_i.
}
\]

In vector notation,

\[
\boxed{
R_p
=
\mathbf{w}^T\mathbf{R}.
}
\]

%%%%%%%%%%%%%%%%%%%%%%%%%%%%%%%%%%%%%%%%%%%%%%%%%%%%%%%%%%%%
\subsection{Expected Portfolio Return}
\label{subsec:expected-portfolio-return}
%%%%%%%%%%%%%%%%%%%%%%%%%%%%%%%%%%%%%%%%%%%%%%%%%%%%%%%%%%%%

Let

\[
\boldsymbol{\mu}
=
\mathbb{E}[\mathbf{R}].
\]

Then

\[
\boxed{
\mathbb{E}[R_p]
=
\mathbf{w}^T\boldsymbol{\mu}.
}
\]

Expected portfolio return is therefore linear in the portfolio weights.

%%%%%%%%%%%%%%%%%%%%%%%%%%%%%%%%%%%%%%%%%%%%%%%%%%%%%%%%%%%%
\subsection{Portfolio Variance}
\label{subsec:portfolio-variance}
%%%%%%%%%%%%%%%%%%%%%%%%%%%%%%%%%%%%%%%%%%%%%%%%%%%%%%%%%%%%

Let covariance matrix

\[
\boxed{
\Sigma
=
\operatorname{Cov}(\mathbf{R}).
}
\]

Then portfolio variance is

\[
\boxed{
\sigma_p^2
=
\mathbf{w}^T
\Sigma
\mathbf{w}.
}
\]

Portfolio standard deviation is

\[
\boxed{
\sigma_p
=
\sqrt{
\mathbf{w}^T\Sigma\mathbf{w}
}.
}
\]

%%%%%%%%%%%%%%%%%%%%%%%%%%%%%%%%%%%%%%%%%%%%%%%%%%%%%%%%%%%%
\subsection{Two-Asset Portfolio Variance}
\label{subsec:two-asset-variance}
%%%%%%%%%%%%%%%%%%%%%%%%%%%%%%%%%%%%%%%%%%%%%%%%%%%%%%%%%%%%

For weights \(w_1,w_2\),

\[
\boxed{
\sigma_p^2
=
w_1^2\sigma_1^2
+
w_2^2\sigma_2^2
+
2w_1w_2
\operatorname{Cov}(R_1,R_2).
}
\]

Using correlation \(\rho_{12}\),

\[
\boxed{
\operatorname{Cov}(R_1,R_2)
=
\rho_{12}\sigma_1\sigma_2.
}
\]

%%%%%%%%%%%%%%%%%%%%%%%%%%%%%%%%%%%%%%%%%%%%%%%%%%%%%%%%%%%%
\subsection{Diversification}
\label{subsec:financial-diversification}
%%%%%%%%%%%%%%%%%%%%%%%%%%%%%%%%%%%%%%%%%%%%%%%%%%%%%%%%%%%%

Diversification reduces portfolio risk when asset returns are not perfectly
positively correlated.

If

\[
\rho_{12}<1,
\]

portfolio variance can be less than a weighted average of individual
variances.

Diversification depends on covariance, not merely on the number of assets.

%%%%%%%%%%%%%%%%%%%%%%%%%%%%%%%%%%%%%%%%%%%%%%%%%%%%%%%%%%%%
\subsection{Worked Example: Diversification}
\label{subsec:worked-diversification}
%%%%%%%%%%%%%%%%%%%%%%%%%%%%%%%%%%%%%%%%%%%%%%%%%%%%%%%%%%%%

\begin{workedexamplebox}{Two Equal-Weight Assets}

Suppose

\[
w_1=w_2=\frac12,
\]

\[
\sigma_1=\sigma_2=0.20,
\]

and

\[
\rho_{12}=0.
\]

Then

\[
\sigma_p^2
=
\frac14(0.20)^2
+
\frac14(0.20)^2.
\]

Thus

\[
\sigma_p^2
=
0.02,
\]

so

\[
\sigma_p
=
\sqrt{0.02}
\approx0.1414.
\]

\begin{answerbox}
\[
\boxed{
\sigma_p\approx14.14\%.
}
\]
\end{answerbox}

The portfolio volatility is lower than either individual volatility.

\end{workedexamplebox}

%%%%%%%%%%%%%%%%%%%%%%%%%%%%%%%%%%%%%%%%%%%%%%%%%%%%%%%%%%%%
\subsection{Mean-Variance Optimization}
\label{subsec:mean-variance-optimization}
%%%%%%%%%%%%%%%%%%%%%%%%%%%%%%%%%%%%%%%%%%%%%%%%%%%%%%%%%%%%

A classical portfolio problem minimizes variance for a target expected
return.

\[
\boxed{
\min_{\mathbf{w}}
\quad
\mathbf{w}^T\Sigma\mathbf{w}
}
\]

subject to

\[
\boxed{
\mathbf{w}^T\boldsymbol{\mu}
=
\mu_p,
}
\]

and

\[
\boxed{
\mathbf{1}^T\mathbf{w}
=
1.
}
\]

Additional constraints may be imposed.

%%%%%%%%%%%%%%%%%%%%%%%%%%%%%%%%%%%%%%%%%%%%%%%%%%%%%%%%%%%%
\subsection{Global Minimum-Variance Portfolio}
\label{subsec:global-minimum-variance}
%%%%%%%%%%%%%%%%%%%%%%%%%%%%%%%%%%%%%%%%%%%%%%%%%%%%%%%%%%%%

Without short-sale constraints, the minimum-variance fully invested portfolio
satisfies

\[
\boxed{
\mathbf{w}_{\mathrm{GMV}}
=
\frac{
\Sigma^{-1}\mathbf{1}
}{
\mathbf{1}^T
\Sigma^{-1}
\mathbf{1}
},
}
\]

assuming \(\Sigma\) is invertible.

This portfolio minimizes variance regardless of target return.

%%%%%%%%%%%%%%%%%%%%%%%%%%%%%%%%%%%%%%%%%%%%%%%%%%%%%%%%%%%%
\subsection{Efficient Frontier}
\label{subsec:efficient-frontier}
%%%%%%%%%%%%%%%%%%%%%%%%%%%%%%%%%%%%%%%%%%%%%%%%%%%%%%%%%%%%

The efficient frontier consists of portfolios providing the greatest
expected return for a given risk level, or equivalently the smallest risk
for a given expected return.

Portfolios below the efficient frontier are dominated.

Thus

\[
\boxed{
\text{efficient portfolio}
=
\text{best attainable risk-return tradeoff}.
}
\]

%%%%%%%%%%%%%%%%%%%%%%%%%%%%%%%%%%%%%%%%%%%%%%%%%%%%%%%%%%%%
\subsection{Risk-Free Asset}
\label{subsec:risk-free-asset}
%%%%%%%%%%%%%%%%%%%%%%%%%%%%%%%%%%%%%%%%%%%%%%%%%%%%%%%%%%%%

A risk-free asset has deterministic return \(r_f\) over the modeled horizon.

Combining a risk-free asset with a risky portfolio produces

\[
\boxed{
R
=
(1-w)r_f+wR_p.
}
\]

Its expected return is linear in \(w\), and its volatility is

\[
\boxed{
\sigma
=
|w|\sigma_p.
}
\]

%%%%%%%%%%%%%%%%%%%%%%%%%%%%%%%%%%%%%%%%%%%%%%%%%%%%%%%%%%%%
\subsection{Sharpe Ratio}
\label{subsec:sharpe-ratio}
%%%%%%%%%%%%%%%%%%%%%%%%%%%%%%%%%%%%%%%%%%%%%%%%%%%%%%%%%%%%

The Sharpe ratio is

\[
\boxed{
\operatorname{SR}
=
\frac{
\mathbb{E}[R_p]-r_f
}{
\sigma_p
}.
}
\]

It measures expected excess return per unit of volatility.

Comparisons require consistent return horizons and assumptions.

%%%%%%%%%%%%%%%%%%%%%%%%%%%%%%%%%%%%%%%%%%%%%%%%%%%%%%%%%%%%
\subsection{Factor Models}
\label{subsec:financial-factor-models}
%%%%%%%%%%%%%%%%%%%%%%%%%%%%%%%%%%%%%%%%%%%%%%%%%%%%%%%%%%%%

A factor model represents asset returns as

\[
\boxed{
R_i
=
\alpha_i
+
\boldsymbol{\beta}_i^T
\mathbf{F}
+
\varepsilon_i.
}
\]

Here:

\[
\mathbf{F}
=
\text{common factors},
\]

\[
\boldsymbol{\beta}_i
=
\text{factor exposures},
\]

and

\[
\varepsilon_i
=
\text{idiosyncratic component}.
\]

Factor models reduce high-dimensional covariance structure.

%%%%%%%%%%%%%%%%%%%%%%%%%%%%%%%%%%%%%%%%%%%%%%%%%%%%%%%%%%%%
\subsection{Market Beta}
\label{subsec:market-beta}
%%%%%%%%%%%%%%%%%%%%%%%%%%%%%%%%%%%%%%%%%%%%%%%%%%%%%%%%%%%%

For asset \(i\) and market return \(R_M\),

\[
\boxed{
\beta_i
=
\frac{
\operatorname{Cov}(R_i,R_M)
}{
\operatorname{Var}(R_M)
}.
}
\]

Beta measures sensitivity to market movements under a linear-factor
interpretation.

%%%%%%%%%%%%%%%%%%%%%%%%%%%%%%%%%%%%%%%%%%%%%%%%%%%%%%%%%%%%
\subsection{Capital Asset Pricing Model}
\label{subsec:capital-asset-pricing-model}
%%%%%%%%%%%%%%%%%%%%%%%%%%%%%%%%%%%%%%%%%%%%%%%%%%%%%%%%%%%%

The classical CAPM relationship is

\[
\boxed{
\mathbb{E}[R_i]
=
r_f
+
\beta_i
\left(
\mathbb{E}[R_M]-r_f
\right).
}
\]

It links expected return to market beta under strong equilibrium assumptions.

CAPM is a model, not a universal empirical law.

%%%%%%%%%%%%%%%%%%%%%%%%%%%%%%%%%%%%%%%%%%%%%%%%%%%%%%%%%%%%
\subsection{Arbitrage}
\label{subsec:financial-arbitrage}
%%%%%%%%%%%%%%%%%%%%%%%%%%%%%%%%%%%%%%%%%%%%%%%%%%%%%%%%%%%%

An arbitrage is a trading strategy that, under the modeled assumptions,
requires no net initial investment, cannot lose, and has positive gain with
positive probability.

Informally,

\[
\boxed{
\text{riskless profit from inconsistent prices}.
}
\]

Absence of arbitrage is a foundational principle of derivative pricing.

%%%%%%%%%%%%%%%%%%%%%%%%%%%%%%%%%%%%%%%%%%%%%%%%%%%%%%%%%%%%
\subsection{Law of One Price}
\label{subsec:law-of-one-price}
%%%%%%%%%%%%%%%%%%%%%%%%%%%%%%%%%%%%%%%%%%%%%%%%%%%%%%%%%%%%

If two portfolios produce exactly the same future cash flows under every
modeled state, absence of arbitrage implies they should have the same current
price.

Thus

\[
\boxed{
\text{same payoff}
\Rightarrow
\text{same price}.
}
\]

This is the law of one price.

%%%%%%%%%%%%%%%%%%%%%%%%%%%%%%%%%%%%%%%%%%%%%%%%%%%%%%%%%%%%
\subsection{Replication}
\label{subsec:financial-replication}
%%%%%%%%%%%%%%%%%%%%%%%%%%%%%%%%%%%%%%%%%%%%%%%%%%%%%%%%%%%%

A derivative is replicated when another portfolio produces the same payoff.

If portfolio value \(V_t\) replicates derivative payoff \(H_T\),

\[
\boxed{
V_T=H_T.
}
\]

Under no-arbitrage assumptions, the derivative price equals the cost of the
replicating portfolio.

%%%%%%%%%%%%%%%%%%%%%%%%%%%%%%%%%%%%%%%%%%%%%%%%%%%%%%%%%%%%
\subsection{Complete Markets}
\label{subsec:complete-markets}
%%%%%%%%%%%%%%%%%%%%%%%%%%%%%%%%%%%%%%%%%%%%%%%%%%%%%%%%%%%%

A market is complete if every relevant contingent claim can be replicated by
traded assets.

In a complete arbitrage-free market, derivative prices are uniquely
determined by replication.

In incomplete markets, several prices may be compatible with absence of
arbitrage.

%%%%%%%%%%%%%%%%%%%%%%%%%%%%%%%%%%%%%%%%%%%%%%%%%%%%%%%%%%%%
\subsection{Derivative Securities}
\label{subsec:derivative-securities}
%%%%%%%%%%%%%%%%%%%%%%%%%%%%%%%%%%%%%%%%%%%%%%%%%%%%%%%%%%%%

A derivative has value determined by another underlying quantity.

Examples include:

\[
\boxed{
\text{forwards},
}
\]

\[
\boxed{
\text{futures},
}
\]

\[
\boxed{
\text{options},
}
\]

and

\[
\boxed{
\text{swaps}.
}
\]

The underlying may be an asset, rate, index, or other measurable quantity.

%%%%%%%%%%%%%%%%%%%%%%%%%%%%%%%%%%%%%%%%%%%%%%%%%%%%%%%%%%%%
\subsection{Forward Contracts}
\label{subsec:forward-contracts}
%%%%%%%%%%%%%%%%%%%%%%%%%%%%%%%%%%%%%%%%%%%%%%%%%%%%%%%%%%%%

A forward contract specifies purchase of an asset at future time \(T\) for
delivery price \(K\).

The long-position payoff is

\[
\boxed{
S_T-K.
}
\]

The short-position payoff is

\[
\boxed{
K-S_T.
}
\]

%%%%%%%%%%%%%%%%%%%%%%%%%%%%%%%%%%%%%%%%%%%%%%%%%%%%%%%%%%%%
\subsection{Forward Price}
\label{subsec:forward-price}
%%%%%%%%%%%%%%%%%%%%%%%%%%%%%%%%%%%%%%%%%%%%%%%%%%%%%%%%%%%%

For a non-dividend-paying asset with current price \(S_0\), constant
continuously compounded risk-free rate \(r\), and no relevant carrying costs,

\[
\boxed{
F_0(T)
=
S_0e^{rT}.
}
\]

Otherwise a cash-and-carry arbitrage could arise under the idealized model.

%%%%%%%%%%%%%%%%%%%%%%%%%%%%%%%%%%%%%%%%%%%%%%%%%%%%%%%%%%%%
\subsection{Worked Example: Forward Price}
\label{subsec:worked-forward-price}
%%%%%%%%%%%%%%%%%%%%%%%%%%%%%%%%%%%%%%%%%%%%%%%%%%%%%%%%%%%%

\begin{workedexamplebox}{Forward on a Non-Dividend Asset}

Suppose

\[
S_0=50,
\qquad
r=0.04,
\qquad
T=1.
\]

Then

\[
F_0
=
50e^{0.04}.
\]

Therefore,

\begin{answerbox}
\[
\boxed{
F_0\approx52.04.
}
\]
\end{answerbox}

\end{workedexamplebox}

%%%%%%%%%%%%%%%%%%%%%%%%%%%%%%%%%%%%%%%%%%%%%%%%%%%%%%%%%%%%
\subsection{European Call Option}
\label{subsec:european-call-option}
%%%%%%%%%%%%%%%%%%%%%%%%%%%%%%%%%%%%%%%%%%%%%%%%%%%%%%%%%%%%

A European call option gives the holder the right, but not the obligation, to
buy the underlying at strike \(K\) at maturity \(T\).

Its payoff is

\[
\boxed{
C_T
=
\max(S_T-K,0).
}
\]

%%%%%%%%%%%%%%%%%%%%%%%%%%%%%%%%%%%%%%%%%%%%%%%%%%%%%%%%%%%%
\subsection{European Put Option}
\label{subsec:european-put-option}
%%%%%%%%%%%%%%%%%%%%%%%%%%%%%%%%%%%%%%%%%%%%%%%%%%%%%%%%%%%%

A European put option gives the holder the right to sell at strike \(K\).

Its payoff is

\[
\boxed{
P_T
=
\max(K-S_T,0).
}
\]

%%%%%%%%%%%%%%%%%%%%%%%%%%%%%%%%%%%%%%%%%%%%%%%%%%%%%%%%%%%%
\subsection{Option Moneyness}
\label{subsec:option-moneyness}
%%%%%%%%%%%%%%%%%%%%%%%%%%%%%%%%%%%%%%%%%%%%%%%%%%%%%%%%%%%%

For a call:

\[
\boxed{
S>K
\Rightarrow
\text{in the money},
}
\]

\[
\boxed{
S=K
\Rightarrow
\text{at the money},
}
\]

and

\[
\boxed{
S<K
\Rightarrow
\text{out of the money}.
}
\]

For puts, the inequalities reverse.

%%%%%%%%%%%%%%%%%%%%%%%%%%%%%%%%%%%%%%%%%%%%%%%%%%%%%%%%%%%%
\subsection{Put-Call Parity}
\label{subsec:put-call-parity}
%%%%%%%%%%%%%%%%%%%%%%%%%%%%%%%%%%%%%%%%%%%%%%%%%%%%%%%%%%%%

For European options on a non-dividend-paying asset under standard
no-arbitrage assumptions,

\[
\boxed{
C_0-P_0
=
S_0-Ke^{-rT}.
}
\]

Equivalently,

\[
\boxed{
C_0+Ke^{-rT}
=
P_0+S_0.
}
\]

%%%%%%%%%%%%%%%%%%%%%%%%%%%%%%%%%%%%%%%%%%%%%%%%%%%%%%%%%%%%
\subsection{Deriving Put-Call Parity}
\label{subsec:deriving-put-call-parity}
%%%%%%%%%%%%%%%%%%%%%%%%%%%%%%%%%%%%%%%%%%%%%%%%%%%%%%%%%%%%

Compare two portfolios:

\[
\boxed{
\text{Portfolio A: call + discounted strike}
}
\]

and

\[
\boxed{
\text{Portfolio B: put + stock}.
}
\]

At maturity, both have value

\[
\boxed{
\max(S_T,K).
}
\]

The law of one price therefore gives put-call parity.

%%%%%%%%%%%%%%%%%%%%%%%%%%%%%%%%%%%%%%%%%%%%%%%%%%%%%%%%%%%%
\subsection{Worked Example: Put-Call Parity}
\label{subsec:worked-put-call-parity}
%%%%%%%%%%%%%%%%%%%%%%%%%%%%%%%%%%%%%%%%%%%%%%%%%%%%%%%%%%%%

\begin{workedexamplebox}{Finding a Put Price}

Suppose

\[
C_0=8,
\quad
S_0=100,
\quad
K=100,
\quad
r=0.05,
\quad
T=1.
\]

From

\[
C_0-P_0
=
S_0-Ke^{-rT},
\]

we obtain

\[
P_0
=
C_0-S_0+Ke^{-rT}.
\]

Thus

\[
P_0
=
8-100+100e^{-0.05}.
\]

Therefore,

\begin{answerbox}
\[
\boxed{
P_0\approx3.12.
}
\]
\end{answerbox}

\end{workedexamplebox}

%%%%%%%%%%%%%%%%%%%%%%%%%%%%%%%%%%%%%%%%%%%%%%%%%%%%%%%%%%%%
\subsection{One-Period Binomial Model}
\label{subsec:one-period-binomial-model}
%%%%%%%%%%%%%%%%%%%%%%%%%%%%%%%%%%%%%%%%%%%%%%%%%%%%%%%%%%%%

Suppose

\[
S_0
\]

can move to either

\[
\boxed{
S_u=uS_0
}
\]

or

\[
\boxed{
S_d=dS_0
}
\]

after one period.

A derivative has payoffs

\[
V_u,
\qquad
V_d.
\]

We seek a portfolio of stock and risk-free borrowing that replicates it.

%%%%%%%%%%%%%%%%%%%%%%%%%%%%%%%%%%%%%%%%%%%%%%%%%%%%%%%%%%%%
\subsection{Binomial Hedge Ratio}
\label{subsec:binomial-hedge-ratio}
%%%%%%%%%%%%%%%%%%%%%%%%%%%%%%%%%%%%%%%%%%%%%%%%%%%%%%%%%%%%

Suppose the replicating portfolio contains \(\Delta\) units of stock and bond
position \(B\).

Then

\[
\Delta S_u+Be^{rT}
=
V_u,
\]

\[
\Delta S_d+Be^{rT}
=
V_d.
\]

Subtracting gives

\[
\boxed{
\Delta
=
\frac{
V_u-V_d
}{
S_u-S_d
}.
}
\]

%%%%%%%%%%%%%%%%%%%%%%%%%%%%%%%%%%%%%%%%%%%%%%%%%%%%%%%%%%%%
\subsection{Risk-Neutral Probability in the Binomial Model}
\label{subsec:binomial-risk-neutral-probability}
%%%%%%%%%%%%%%%%%%%%%%%%%%%%%%%%%%%%%%%%%%%%%%%%%%%%%%%%%%%%

Define

\[
\boxed{
q
=
\frac{
e^{rT}-d
}{
u-d
}.
}
\]

Under the no-arbitrage condition

\[
\boxed{
d<e^{rT}<u,
}
\]

we have

\[
0<q<1.
\]

The derivative price is

\[
\boxed{
V_0
=
e^{-rT}
\left[
qV_u+(1-q)V_d
\right].
}
\]

%%%%%%%%%%%%%%%%%%%%%%%%%%%%%%%%%%%%%%%%%%%%%%%%%%%%%%%%%%%%
\subsection{Worked Example: Binomial Call}
\label{subsec:worked-binomial-call}
%%%%%%%%%%%%%%%%%%%%%%%%%%%%%%%%%%%%%%%%%%%%%%%%%%%%%%%%%%%%

\begin{workedexamplebox}{One-Period Option Pricing}

Suppose

\[
S_0=100,
\qquad
u=1.2,
\qquad
d=0.8,
\]

and assume the one-period risk-free growth factor is

\[
e^{rT}=1.05.
\]

Consider a call with

\[
K=100.
\]

Then

\[
S_u=120,
\qquad
S_d=80.
\]

The payoffs are

\[
V_u=20,
\qquad
V_d=0.
\]

The risk-neutral probability is

\[
q
=
\frac{1.05-0.8}{1.2-0.8}
=
0.625.
\]

Therefore,

\[
V_0
=
\frac{1}{1.05}
\left[
0.625(20)
\right].
\]

Hence

\begin{answerbox}
\[
\boxed{
V_0\approx11.90.
}
\]
\end{answerbox}

\end{workedexamplebox}

%%%%%%%%%%%%%%%%%%%%%%%%%%%%%%%%%%%%%%%%%%%%%%%%%%%%%%%%%%%%
\subsection{Multi-Period Binomial Trees}
\label{subsec:multi-period-binomial-trees}
%%%%%%%%%%%%%%%%%%%%%%%%%%%%%%%%%%%%%%%%%%%%%%%%%%%%%%%%%%%%

A multi-period binomial model generates a price tree.

Derivative values are computed backward from terminal payoffs.

At each node,

\[
\boxed{
V
=
e^{-r\Delta t}
\left[
qV_u+(1-q)V_d
\right].
}
\]

This is backward induction.

%%%%%%%%%%%%%%%%%%%%%%%%%%%%%%%%%%%%%%%%%%%%%%%%%%%%%%%%%%%%
\subsection{American Options}
\label{subsec:american-options}
%%%%%%%%%%%%%%%%%%%%%%%%%%%%%%%%%%%%%%%%%%%%%%%%%%%%%%%%%%%%

An American option may be exercised at any time up to maturity.

At a binomial-tree node,

\[
\boxed{
V
=
\max
\left\{
\text{exercise value},
\text{continuation value}
\right\}.
}
\]

This creates an optimal stopping problem.

%%%%%%%%%%%%%%%%%%%%%%%%%%%%%%%%%%%%%%%%%%%%%%%%%%%%%%%%%%%%
\subsection{Stochastic Asset Prices}
\label{subsec:stochastic-asset-prices}
%%%%%%%%%%%%%%%%%%%%%%%%%%%%%%%%%%%%%%%%%%%%%%%%%%%%%%%%%%%%

Financial prices are uncertain.

A common continuous-time model is

\[
\boxed{
dS_t
=
\mu S_t\,dt
+
\sigma S_t\,dW_t.
}
\]

Here:

\[
\mu
=
\text{drift},
\]

\[
\sigma
=
\text{volatility},
\]

and

\[
W_t
=
\text{Brownian motion}.
\]

%%%%%%%%%%%%%%%%%%%%%%%%%%%%%%%%%%%%%%%%%%%%%%%%%%%%%%%%%%%%
\subsection{Geometric Brownian Motion}
\label{subsec:geometric-brownian-motion-finance}
%%%%%%%%%%%%%%%%%%%%%%%%%%%%%%%%%%%%%%%%%%%%%%%%%%%%%%%%%%%%

The stochastic differential equation

\[
dS_t
=
\mu S_tdt
+
\sigma S_tdW_t
\]

defines geometric Brownian motion.

Its solution is

\[
\boxed{
S_t
=
S_0
\exp
\left[
\left(
\mu-\frac12\sigma^2
\right)t
+
\sigma W_t
\right].
}
\]

Thus \(S_t>0\) whenever \(S_0>0\).

%%%%%%%%%%%%%%%%%%%%%%%%%%%%%%%%%%%%%%%%%%%%%%%%%%%%%%%%%%%%
\subsection{Log Price Dynamics}
\label{subsec:log-price-dynamics}
%%%%%%%%%%%%%%%%%%%%%%%%%%%%%%%%%%%%%%%%%%%%%%%%%%%%%%%%%%%%

Applying It\^o's formula to

\[
\ln S_t
\]

gives

\[
\boxed{
d\ln S_t
=
\left(
\mu-\frac12\sigma^2
\right)dt
+
\sigma dW_t.
}
\]

The correction

\[
-\frac12\sigma^2
\]

is an It\^o-calculus effect.

%%%%%%%%%%%%%%%%%%%%%%%%%%%%%%%%%%%%%%%%%%%%%%%%%%%%%%%%%%%%
\subsection{Volatility}
\label{subsec:financial-volatility}
%%%%%%%%%%%%%%%%%%%%%%%%%%%%%%%%%%%%%%%%%%%%%%%%%%%%%%%%%%%%

Volatility measures the scale of return fluctuations.

In geometric Brownian motion, the parameter

\[
\boxed{
\sigma
}
\]

controls diffusion intensity.

Over a short interval \(\Delta t\),

\[
\boxed{
\operatorname{Var}
(\Delta\ln S)
\approx
\sigma^2\Delta t.
}
\]

%%%%%%%%%%%%%%%%%%%%%%%%%%%%%%%%%%%%%%%%%%%%%%%%%%%%%%%%%%%%
\subsection{Historical Volatility}
\label{subsec:historical-volatility}
%%%%%%%%%%%%%%%%%%%%%%%%%%%%%%%%%%%%%%%%%%%%%%%%%%%%%%%%%%%%

Historical volatility is estimated from observed returns.

For log returns

\[
r_1,\ldots,r_n,
\]

one may estimate standard deviation and annualize it according to the
sampling frequency.

The estimate depends on:

\[
\boxed{
\text{sample period},
}
\]

\[
\boxed{
\text{return frequency},
}
\]

and

\[
\boxed{
\text{estimation method}.
}
\]

%%%%%%%%%%%%%%%%%%%%%%%%%%%%%%%%%%%%%%%%%%%%%%%%%%%%%%%%%%%%
\subsection{Implied Volatility}
\label{subsec:implied-volatility}
%%%%%%%%%%%%%%%%%%%%%%%%%%%%%%%%%%%%%%%%%%%%%%%%%%%%%%%%%%%%

Implied volatility is the volatility parameter that makes an option-pricing
model reproduce an observed market price.

Schematically,

\[
\boxed{
C_{\mathrm{model}}(\sigma_{\mathrm{impl}})
=
C_{\mathrm{market}}.
}
\]

It is obtained by numerical inversion.

%%%%%%%%%%%%%%%%%%%%%%%%%%%%%%%%%%%%%%%%%%%%%%%%%%%%%%%%%%%%
\subsection{Risk-Neutral Valuation}
\label{subsec:risk-neutral-valuation}
%%%%%%%%%%%%%%%%%%%%%%%%%%%%%%%%%%%%%%%%%%%%%%%%%%%%%%%%%%%%

Under an equivalent risk-neutral measure \(Q\), discounted tradable asset
prices are martingales under suitable market assumptions.

Derivative value can then be written

\[
\boxed{
V_t
=
\mathbb{E}^{Q}
\left[
e^{-r(T-t)}
H_T
\mid
\mathcal{F}_t
\right]
}
\]

for constant risk-free rate \(r\).

The expectation uses the risk-neutral probability measure, not necessarily
the real-world probability distribution.

%%%%%%%%%%%%%%%%%%%%%%%%%%%%%%%%%%%%%%%%%%%%%%%%%%%%%%%%%%%%
\subsection{Risk-Neutral Asset Dynamics}
\label{subsec:risk-neutral-asset-dynamics}
%%%%%%%%%%%%%%%%%%%%%%%%%%%%%%%%%%%%%%%%%%%%%%%%%%%%%%%%%%%%

Under the standard Black--Scholes assumptions,

\[
\boxed{
dS_t
=
rS_t\,dt
+
\sigma S_t\,dW_t^{Q}.
}
\]

The real-world drift \(\mu\) is replaced by the risk-free rate \(r\) for
pricing.

This arises from replication and no-arbitrage rather than from assuming
investors believe the true expected return equals \(r\).

%%%%%%%%%%%%%%%%%%%%%%%%%%%%%%%%%%%%%%%%%%%%%%%%%%%%%%%%%%%%
\subsection{Black--Scholes Model}
\label{subsec:black-scholes-model}
%%%%%%%%%%%%%%%%%%%%%%%%%%%%%%%%%%%%%%%%%%%%%%%%%%%%%%%%%%%%

The Black--Scholes framework assumes, in its basic form:

\begin{itemize}
    \item geometric Brownian motion for the underlying;
    \item constant volatility;
    \item constant risk-free rate;
    \item continuous trading;
    \item frictionless markets;
    \item no arbitrage;
    \item unrestricted borrowing and short selling in the idealized model.
\end{itemize}

These assumptions produce a tractable derivative-pricing framework.

%%%%%%%%%%%%%%%%%%%%%%%%%%%%%%%%%%%%%%%%%%%%%%%%%%%%%%%%%%%%
\subsection{Black--Scholes PDE}
\label{subsec:black-scholes-pde}
%%%%%%%%%%%%%%%%%%%%%%%%%%%%%%%%%%%%%%%%%%%%%%%%%%%%%%%%%%%%

Let derivative value be

\[
V=V(S,t).
\]

Using It\^o's formula and a locally riskless hedged portfolio leads to

\[
\boxed{
\frac{\partial V}{\partial t}
+
\frac12\sigma^2S^2
\frac{\partial^2V}{\partial S^2}
+
rS
\frac{\partial V}{\partial S}
-
rV
=
0.
}
\]

This is the Black--Scholes partial differential equation.

%%%%%%%%%%%%%%%%%%%%%%%%%%%%%%%%%%%%%%%%%%%%%%%%%%%%%%%%%%%%
\subsection{European Call Formula}
\label{subsec:black-scholes-call-formula}
%%%%%%%%%%%%%%%%%%%%%%%%%%%%%%%%%%%%%%%%%%%%%%%%%%%%%%%%%%%%

For a non-dividend-paying stock, the Black--Scholes European call price is

\[
\boxed{
C
=
S_0\Phi(d_1)
-
Ke^{-rT}\Phi(d_2),
}
\]

where

\[
\boxed{
d_1
=
\frac{
\ln(S_0/K)
+
(r+\frac12\sigma^2)T
}{
\sigma\sqrt{T}
},
}
\]

and

\[
\boxed{
d_2
=
d_1-\sigma\sqrt{T}.
}
\]

Here \(\Phi\) is the standard normal cumulative distribution function.

%%%%%%%%%%%%%%%%%%%%%%%%%%%%%%%%%%%%%%%%%%%%%%%%%%%%%%%%%%%%
\subsection{European Put Formula}
\label{subsec:black-scholes-put-formula}
%%%%%%%%%%%%%%%%%%%%%%%%%%%%%%%%%%%%%%%%%%%%%%%%%%%%%%%%%%%%

The corresponding European put price is

\[
\boxed{
P
=
Ke^{-rT}
\Phi(-d_2)
-
S_0\Phi(-d_1).
}
\]

This is consistent with put-call parity.

%%%%%%%%%%%%%%%%%%%%%%%%%%%%%%%%%%%%%%%%%%%%%%%%%%%%%%%%%%%%
\subsection{Delta}
\label{subsec:option-delta}
%%%%%%%%%%%%%%%%%%%%%%%%%%%%%%%%%%%%%%%%%%%%%%%%%%%%%%%%%%%%

Delta measures sensitivity to the underlying price:

\[
\boxed{
\Delta
=
\frac{\partial V}{\partial S}.
}
\]

For a Black--Scholes call,

\[
\boxed{
\Delta_C
=
\Phi(d_1).
}
\]

Delta is central to local hedging.

%%%%%%%%%%%%%%%%%%%%%%%%%%%%%%%%%%%%%%%%%%%%%%%%%%%%%%%%%%%%
\subsection{Gamma}
\label{subsec:option-gamma}
%%%%%%%%%%%%%%%%%%%%%%%%%%%%%%%%%%%%%%%%%%%%%%%%%%%%%%%%%%%%

Gamma measures sensitivity of delta to the underlying:

\[
\boxed{
\Gamma
=
\frac{\partial^2V}{\partial S^2}.
}
\]

Large gamma means delta changes rapidly when the underlying moves.

Gamma therefore measures curvature risk.

%%%%%%%%%%%%%%%%%%%%%%%%%%%%%%%%%%%%%%%%%%%%%%%%%%%%%%%%%%%%
\subsection{Theta}
\label{subsec:option-theta}
%%%%%%%%%%%%%%%%%%%%%%%%%%%%%%%%%%%%%%%%%%%%%%%%%%%%%%%%%%%%

Theta is

\[
\boxed{
\Theta
=
\frac{\partial V}{\partial t}.
}
\]

It measures sensitivity to passage of time, holding other model inputs fixed.

Option-value decay is often associated with theta, though its sign depends on
the instrument and conventions.

%%%%%%%%%%%%%%%%%%%%%%%%%%%%%%%%%%%%%%%%%%%%%%%%%%%%%%%%%%%%
\subsection{Vega}
\label{subsec:option-vega}
%%%%%%%%%%%%%%%%%%%%%%%%%%%%%%%%%%%%%%%%%%%%%%%%%%%%%%%%%%%%

Vega measures sensitivity to volatility:

\[
\boxed{
\operatorname{Vega}
=
\frac{\partial V}{\partial\sigma}.
}
\]

Options can be highly sensitive to volatility even when the underlying price
does not change.

%%%%%%%%%%%%%%%%%%%%%%%%%%%%%%%%%%%%%%%%%%%%%%%%%%%%%%%%%%%%
\subsection{Rho}
\label{subsec:option-rho}
%%%%%%%%%%%%%%%%%%%%%%%%%%%%%%%%%%%%%%%%%%%%%%%%%%%%%%%%%%%%

Rho measures sensitivity to the risk-free interest rate:

\[
\boxed{
\rho_V
=
\frac{\partial V}{\partial r}.
}
\]

Its importance generally increases for longer-maturity derivatives.

%%%%%%%%%%%%%%%%%%%%%%%%%%%%%%%%%%%%%%%%%%%%%%%%%%%%%%%%%%%%
\subsection{Delta Hedging}
\label{subsec:delta-hedging}
%%%%%%%%%%%%%%%%%%%%%%%%%%%%%%%%%%%%%%%%%%%%%%%%%%%%%%%%%%%%

Suppose a derivative has delta \(\Delta\).

A locally delta-neutral portfolio can be constructed by offsetting the
underlying exposure.

For one derivative position,

\[
\boxed{
\Pi
=
V-\Delta S.
}
\]

Over an infinitesimal interval in the Black--Scholes model, the stochastic
first-order price exposure can be removed by appropriate delta hedging.

%%%%%%%%%%%%%%%%%%%%%%%%%%%%%%%%%%%%%%%%%%%%%%%%%%%%%%%%%%%%
\subsection{Dynamic Hedging}
\label{subsec:dynamic-hedging}
%%%%%%%%%%%%%%%%%%%%%%%%%%%%%%%%%%%%%%%%%%%%%%%%%%%%%%%%%%%%

Delta changes as:

\[
\boxed{
S,
}
\]

\[
\boxed{
t,
}
\]

and

\[
\boxed{
\sigma
}
\]

change.

Therefore a delta hedge generally requires repeated rebalancing.

In real markets, discrete trading and transaction costs prevent perfect
continuous replication.

%%%%%%%%%%%%%%%%%%%%%%%%%%%%%%%%%%%%%%%%%%%%%%%%%%%%%%%%%%%%
\subsection{Volatility Smile}
\label{subsec:volatility-smile}
%%%%%%%%%%%%%%%%%%%%%%%%%%%%%%%%%%%%%%%%%%%%%%%%%%%%%%%%%%%%

If Black--Scholes assumptions were exact with constant volatility, options
with the same maturity would imply the same \(\sigma\).

In practice, implied volatility often varies with strike.

This pattern is called a

\[
\boxed{
\text{volatility smile}
}
\]

or

\[
\boxed{
\text{volatility skew}.
}
\]

%%%%%%%%%%%%%%%%%%%%%%%%%%%%%%%%%%%%%%%%%%%%%%%%%%%%%%%%%%%%
\subsection{Volatility Surface}
\label{subsec:volatility-surface}
%%%%%%%%%%%%%%%%%%%%%%%%%%%%%%%%%%%%%%%%%%%%%%%%%%%%%%%%%%%%

Implied volatility depends on strike and maturity:

\[
\boxed{
\sigma_{\mathrm{impl}}
=
\sigma_{\mathrm{impl}}(K,T).
}
\]

The resulting surface summarizes market option prices in volatility terms.

It also demonstrates limitations of constant-volatility models.

%%%%%%%%%%%%%%%%%%%%%%%%%%%%%%%%%%%%%%%%%%%%%%%%%%%%%%%%%%%%
\subsection{Monte Carlo Derivative Pricing}
\label{subsec:monte-carlo-derivative-pricing}
%%%%%%%%%%%%%%%%%%%%%%%%%%%%%%%%%%%%%%%%%%%%%%%%%%%%%%%%%%%%

Under risk-neutral valuation,

\[
V_0
=
e^{-rT}
\mathbb{E}^{Q}[H_T].
\]

Simulate independent payoffs

\[
H_T^{(1)},
\ldots,
H_T^{(N)}.
\]

Then estimate

\[
\boxed{
\hat V_0
=
e^{-rT}
\frac1N
\sum_{i=1}^{N}
H_T^{(i)}.
}
\]

%%%%%%%%%%%%%%%%%%%%%%%%%%%%%%%%%%%%%%%%%%%%%%%%%%%%%%%%%%%%
\subsection{Monte Carlo Error}
\label{subsec:monte-carlo-finance-error}
%%%%%%%%%%%%%%%%%%%%%%%%%%%%%%%%%%%%%%%%%%%%%%%%%%%%%%%%%%%%

Under standard finite-variance conditions, Monte Carlo standard error scales
as

\[
\boxed{
O(N^{-1/2}).
}
\]

Thus reducing error by a factor of \(10\) typically requires approximately
\(100\) times as many independent simulations.

%%%%%%%%%%%%%%%%%%%%%%%%%%%%%%%%%%%%%%%%%%%%%%%%%%%%%%%%%%%%
\subsection{Variance Reduction}
\label{subsec:finance-variance-reduction}
%%%%%%%%%%%%%%%%%%%%%%%%%%%%%%%%%%%%%%%%%%%%%%%%%%%%%%%%%%%%

Variance-reduction techniques seek greater accuracy without proportionally
increasing simulation count.

Examples include:

\[
\boxed{
\text{antithetic variates},
}
\]

\[
\boxed{
\text{control variates},
}
\]

\[
\boxed{
\text{importance sampling},
}
\]

and

\[
\boxed{
\text{stratification}.
}
\]

%%%%%%%%%%%%%%%%%%%%%%%%%%%%%%%%%%%%%%%%%%%%%%%%%%%%%%%%%%%%
\subsection{Finite-Difference Option Pricing}
\label{subsec:finite-difference-option-pricing}
%%%%%%%%%%%%%%%%%%%%%%%%%%%%%%%%%%%%%%%%%%%%%%%%%%%%%%%%%%%%

The Black--Scholes PDE can be discretized in asset price and time.

For a grid

\[
S_i=i\Delta S,
\qquad
t_n=n\Delta t,
\]

finite differences approximate

\[
V_t,
\qquad
V_S,
\qquad
V_{SS}.
\]

The terminal payoff provides the final-time condition.

%%%%%%%%%%%%%%%%%%%%%%%%%%%%%%%%%%%%%%%%%%%%%%%%%%%%%%%%%%%%
\subsection{Boundary Conditions in Option Pricing}
\label{subsec:option-pricing-boundary-conditions}
%%%%%%%%%%%%%%%%%%%%%%%%%%%%%%%%%%%%%%%%%%%%%%%%%%%%%%%%%%%%

For a European call,

\[
\boxed{
V(S,T)
=
\max(S-K,0).
}
\]

Additional conditions are imposed as \(S\) approaches the numerical
boundaries.

For example,

\[
\boxed{
V(0,t)=0
}
\]

for a standard call on a nonnegative underlying.

%%%%%%%%%%%%%%%%%%%%%%%%%%%%%%%%%%%%%%%%%%%%%%%%%%%%%%%%%%%%
\subsection{Stochastic Volatility}
\label{subsec:stochastic-volatility}
%%%%%%%%%%%%%%%%%%%%%%%%%%%%%%%%%%%%%%%%%%%%%%%%%%%%%%%%%%%%

Constant volatility is often unrealistic.

A stochastic-volatility model lets variance itself evolve randomly:

\[
\boxed{
dS_t
=
\mu S_tdt
+
\sqrt{v_t}S_tdW_t^{(1)},
}
\]

together with a stochastic equation for

\[
\boxed{
v_t.
}
\]

Such models can better represent changing volatility and implied-volatility
patterns.

%%%%%%%%%%%%%%%%%%%%%%%%%%%%%%%%%%%%%%%%%%%%%%%%%%%%%%%%%%%%
\subsection{Jump-Diffusion Models}
\label{subsec:jump-diffusion-models}
%%%%%%%%%%%%%%%%%%%%%%%%%%%%%%%%%%%%%%%%%%%%%%%%%%%%%%%%%%%%

Continuous Brownian models cannot directly generate instantaneous price
jumps.

A jump-diffusion model adds a jump process:

\[
\boxed{
\frac{dS_t}{S_{t^-}}
=
\mu\,dt
+
\sigma\,dW_t
+
dJ_t.
}
\]

The process \(J_t\) represents discontinuous price changes.

%%%%%%%%%%%%%%%%%%%%%%%%%%%%%%%%%%%%%%%%%%%%%%%%%%%%%%%%%%%%
\subsection{Interest-Rate Term Structure}
\label{subsec:interest-rate-term-structure}
%%%%%%%%%%%%%%%%%%%%%%%%%%%%%%%%%%%%%%%%%%%%%%%%%%%%%%%%%%%%

Interest rates depend on maturity.

The term structure describes rates across different horizons.

Common representations include:

\[
\boxed{
\text{discount factors},
}
\]

\[
\boxed{
\text{spot rates},
}
\]

\[
\boxed{
\text{forward rates},
}
\]

and

\[
\boxed{
\text{yield curves}.
}
\]

%%%%%%%%%%%%%%%%%%%%%%%%%%%%%%%%%%%%%%%%%%%%%%%%%%%%%%%%%%%%
\subsection{Spot Rates}
\label{subsec:spot-rates}
%%%%%%%%%%%%%%%%%%%%%%%%%%%%%%%%%%%%%%%%%%%%%%%%%%%%%%%%%%%%

A continuously compounded spot rate \(R(0,T)\) satisfies

\[
\boxed{
P(0,T)
=
e^{-R(0,T)T}
}
\]

for a zero-coupon bond paying one unit at maturity.

Thus

\[
\boxed{
R(0,T)
=
-\frac1T
\ln P(0,T).
}
\]

%%%%%%%%%%%%%%%%%%%%%%%%%%%%%%%%%%%%%%%%%%%%%%%%%%%%%%%%%%%%
\subsection{Forward Rates}
\label{subsec:forward-rates}
%%%%%%%%%%%%%%%%%%%%%%%%%%%%%%%%%%%%%%%%%%%%%%%%%%%%%%%%%%%%

A forward rate describes a future borrowing or investment rate implied by
current discount factors.

Under continuous compounding, the instantaneous forward rate can be written

\[
\boxed{
f(0,T)
=
-
\frac{\partial}{\partial T}
\ln P(0,T).
}
\]

This connects discount curves with future implied rates.

%%%%%%%%%%%%%%%%%%%%%%%%%%%%%%%%%%%%%%%%%%%%%%%%%%%%%%%%%%%%
\subsection{Short-Rate Models}
\label{subsec:short-rate-models}
%%%%%%%%%%%%%%%%%%%%%%%%%%%%%%%%%%%%%%%%%%%%%%%%%%%%%%%%%%%%

An interest-rate model may specify dynamics of the instantaneous short rate

\[
\boxed{
r_t.
}
\]

A generic stochastic model is

\[
\boxed{
dr_t
=
a(r_t,t)\,dt
+
b(r_t,t)\,dW_t.
}
\]

Bond prices are then derived from the modeled rate process.

%%%%%%%%%%%%%%%%%%%%%%%%%%%%%%%%%%%%%%%%%%%%%%%%%%%%%%%%%%%%
\subsection{Mean-Reverting Interest Rates}
\label{subsec:mean-reverting-interest-rates}
%%%%%%%%%%%%%%%%%%%%%%%%%%%%%%%%%%%%%%%%%%%%%%%%%%%%%%%%%%%%

A simple mean-reverting structure is

\[
\boxed{
dr_t
=
\kappa(\theta-r_t)\,dt
+
\sigma\,dW_t.
}
\]

The drift pulls \(r_t\) toward long-run level \(\theta\).

The parameter \(\kappa\) controls the speed of mean reversion.

%%%%%%%%%%%%%%%%%%%%%%%%%%%%%%%%%%%%%%%%%%%%%%%%%%%%%%%%%%%%
\subsection{Credit Risk}
\label{subsec:credit-risk}
%%%%%%%%%%%%%%%%%%%%%%%%%%%%%%%%%%%%%%%%%%%%%%%%%%%%%%%%%%%%

Credit risk is the possibility that a borrower or counterparty fails to make
promised payments.

Basic components include:

\[
\boxed{
\text{probability of default},
}
\]

\[
\boxed{
\text{exposure at default},
}
\]

and

\[
\boxed{
\text{loss given default}.
}
\]

%%%%%%%%%%%%%%%%%%%%%%%%%%%%%%%%%%%%%%%%%%%%%%%%%%%%%%%%%%%%
\subsection{Expected Credit Loss}
\label{subsec:expected-credit-loss}
%%%%%%%%%%%%%%%%%%%%%%%%%%%%%%%%%%%%%%%%%%%%%%%%%%%%%%%%%%%%

A basic approximation is

\[
\boxed{
\operatorname{ECL}
=
\operatorname{PD}
\times
\operatorname{LGD}
\times
\operatorname{EAD},
}
\]

where

\[
\operatorname{PD}
=
\text{probability of default},
\]

\[
\operatorname{LGD}
=
\text{loss given default},
\]

and

\[
\operatorname{EAD}
=
\text{exposure at default}.
\]

%%%%%%%%%%%%%%%%%%%%%%%%%%%%%%%%%%%%%%%%%%%%%%%%%%%%%%%%%%%%
\subsection{Market Risk}
\label{subsec:market-risk}
%%%%%%%%%%%%%%%%%%%%%%%%%%%%%%%%%%%%%%%%%%%%%%%%%%%%%%%%%%%%

Market risk arises from movements in quantities such as:

\[
\boxed{
\text{asset prices},
}
\]

\[
\boxed{
\text{interest rates},
}
\]

\[
\boxed{
\text{exchange rates},
}
\]

and

\[
\boxed{
\text{volatility}.
}
\]

Risk measurement attempts to summarize potential portfolio losses.

%%%%%%%%%%%%%%%%%%%%%%%%%%%%%%%%%%%%%%%%%%%%%%%%%%%%%%%%%%%%
\subsection{Value at Risk}
\label{subsec:value-at-risk-finance}
%%%%%%%%%%%%%%%%%%%%%%%%%%%%%%%%%%%%%%%%%%%%%%%%%%%%%%%%%%%%

Let \(L\) denote loss.

For confidence level \(\alpha\), Value at Risk can be defined as a loss
quantile:

\[
\boxed{
\operatorname{VaR}_{\alpha}(L)
=
\inf
\left\{
\ell:
\Pr(L\leq\ell)\geq\alpha
\right\}.
}
\]

It describes a threshold exceeded only with probability approximately
\(1-\alpha\).

%%%%%%%%%%%%%%%%%%%%%%%%%%%%%%%%%%%%%%%%%%%%%%%%%%%%%%%%%%%%
\subsection{VaR Interpretation}
\label{subsec:var-interpretation-finance}
%%%%%%%%%%%%%%%%%%%%%%%%%%%%%%%%%%%%%%%%%%%%%%%%%%%%%%%%%%%%

If

\[
\operatorname{VaR}_{0.99}
=
\$1{,}000{,}000
\]

over one day, the model says that approximately \(99\%\) of one-day losses
do not exceed one million dollars.

It does not specify how large losses can be beyond that threshold.

%%%%%%%%%%%%%%%%%%%%%%%%%%%%%%%%%%%%%%%%%%%%%%%%%%%%%%%%%%%%
\subsection{Conditional Value at Risk}
\label{subsec:cvar-finance}
%%%%%%%%%%%%%%%%%%%%%%%%%%%%%%%%%%%%%%%%%%%%%%%%%%%%%%%%%%%%

Conditional Value at Risk focuses on tail losses beyond a quantile.

A useful optimization representation is

\[
\boxed{
\operatorname{CVaR}_{\alpha}(L)
=
\min_{\eta}
\left[
\eta
+
\frac{1}{1-\alpha}
\mathbb{E}(L-\eta)_+
\right],
}
\]

where

\[
(x)_+
=
\max(x,0).
\]

%%%%%%%%%%%%%%%%%%%%%%%%%%%%%%%%%%%%%%%%%%%%%%%%%%%%%%%%%%%%
\subsection{VaR Versus CVaR}
\label{subsec:var-versus-cvar}
%%%%%%%%%%%%%%%%%%%%%%%%%%%%%%%%%%%%%%%%%%%%%%%%%%%%%%%%%%%%

VaR describes a quantile threshold.

CVaR incorporates the severity of losses in the tail.

Thus

\[
\boxed{
\text{VaR}
\rightarrow
\text{where the tail begins},
}
\]

while

\[
\boxed{
\text{CVaR}
\rightarrow
\text{how severe the tail is}.
}
\]

%%%%%%%%%%%%%%%%%%%%%%%%%%%%%%%%%%%%%%%%%%%%%%%%%%%%%%%%%%%%
\subsection{Risk-Constrained Portfolio Optimization}
\label{subsec:risk-constrained-portfolio}
%%%%%%%%%%%%%%%%%%%%%%%%%%%%%%%%%%%%%%%%%%%%%%%%%%%%%%%%%%%%

A portfolio may maximize expected return while imposing a risk limit:

\[
\boxed{
\max_{\mathbf{w}}
\quad
\mathbf{w}^T\boldsymbol{\mu}
}
\]

subject to

\[
\boxed{
\operatorname{CVaR}_{\alpha}
(L(\mathbf{w}))
\leq
c.
}
\]

This connects mathematical finance directly with convex and stochastic
optimization.

%%%%%%%%%%%%%%%%%%%%%%%%%%%%%%%%%%%%%%%%%%%%%%%%%%%%%%%%%%%%
\subsection{Dynamic Portfolio Choice}
\label{subsec:dynamic-portfolio-choice}
%%%%%%%%%%%%%%%%%%%%%%%%%%%%%%%%%%%%%%%%%%%%%%%%%%%%%%%%%%%%

Suppose wealth evolves as

\[
\boxed{
dW_t
=
\text{investment return}
+
\text{cash flows}.
}
\]

An investor chooses portfolio weights over time.

The optimization problem becomes a stochastic control problem.

The policy may depend on:

\[
\boxed{
W_t,
}
\]

\[
\boxed{
t,
}
\]

and other state variables.

%%%%%%%%%%%%%%%%%%%%%%%%%%%%%%%%%%%%%%%%%%%%%%%%%%%%%%%%%%%%
\subsection{Transaction Costs}
\label{subsec:transaction-costs-finance}
%%%%%%%%%%%%%%%%%%%%%%%%%%%%%%%%%%%%%%%%%%%%%%%%%%%%%%%%%%%%

Real trading incurs costs.

A simplified proportional transaction cost may be

\[
\boxed{
c|\Delta q|,
}
\]

where \(\Delta q\) is the change in position.

Transaction costs can make continuous rebalancing economically impossible.

They also alter optimal trading and hedging policies.

%%%%%%%%%%%%%%%%%%%%%%%%%%%%%%%%%%%%%%%%%%%%%%%%%%%%%%%%%%%%
\subsection{Bid-Ask Spread}
\label{subsec:bid-ask-spread}
%%%%%%%%%%%%%%%%%%%%%%%%%%%%%%%%%%%%%%%%%%%%%%%%%%%%%%%%%%%%

A market may quote:

\[
\boxed{
S_{\mathrm{bid}}
}
\]

for the price at which one can sell, and

\[
\boxed{
S_{\mathrm{ask}}
}
\]

for the price at which one can buy.

The spread is

\[
\boxed{
S_{\mathrm{ask}}
-
S_{\mathrm{bid}}.
}
\]

This is an important source of trading friction.

%%%%%%%%%%%%%%%%%%%%%%%%%%%%%%%%%%%%%%%%%%%%%%%%%%%%%%%%%%%%
\subsection{Liquidity Risk}
\label{subsec:liquidity-risk}
%%%%%%%%%%%%%%%%%%%%%%%%%%%%%%%%%%%%%%%%%%%%%%%%%%%%%%%%%%%%

Liquidity risk arises when positions cannot be traded quickly at expected
prices.

A large order may move the market.

Thus the realized execution price can depend on:

\[
\boxed{
\text{trade size},
}
\]

\[
\boxed{
\text{market depth},
}
\]

and

\[
\boxed{
\text{trading speed}.
}
\]

%%%%%%%%%%%%%%%%%%%%%%%%%%%%%%%%%%%%%%%%%%%%%%%%%%%%%%%%%%%%
\subsection{Market Impact}
\label{subsec:market-impact}
%%%%%%%%%%%%%%%%%%%%%%%%%%%%%%%%%%%%%%%%%%%%%%%%%%%%%%%%%%%%

Market impact means a trade itself changes available prices.

A simple model might write execution cost as

\[
\boxed{
C(q)
=
S_0q
+
\lambda q^2,
}
\]

where the quadratic term represents increasing impact with order size.

This turns trading into an optimization problem with endogenous cost.

%%%%%%%%%%%%%%%%%%%%%%%%%%%%%%%%%%%%%%%%%%%%%%%%%%%%%%%%%%%%
\subsection{Calibration}
\label{subsec:financial-calibration}
%%%%%%%%%%%%%%%%%%%%%%%%%%%%%%%%%%%%%%%%%%%%%%%%%%%%%%%%%%%%

Calibration chooses model parameters to match observed prices.

Suppose market prices are

\[
P_1^{\mathrm{mkt}},
\ldots,
P_n^{\mathrm{mkt}}.
\]

A least-squares calibration may solve

\[
\boxed{
\min_{\boldsymbol{\theta}}
\sum_{i=1}^{n}
\left[
P_i^{\mathrm{model}}(\boldsymbol{\theta})
-
P_i^{\mathrm{mkt}}
\right]^2.
}
\]

%%%%%%%%%%%%%%%%%%%%%%%%%%%%%%%%%%%%%%%%%%%%%%%%%%%%%%%%%%%%
\subsection{Parameter Estimation}
\label{subsec:financial-parameter-estimation}
%%%%%%%%%%%%%%%%%%%%%%%%%%%%%%%%%%%%%%%%%%%%%%%%%%%%%%%%%%%%

Historical data may be used to estimate:

\[
\boxed{
\text{drift},
}
\]

\[
\boxed{
\text{volatility},
}
\]

\[
\boxed{
\text{correlations},
}
\]

\[
\boxed{
\text{factor exposures},
}
\]

and

\[
\boxed{
\text{default probabilities}.
}
\]

These estimates contain sampling uncertainty.

%%%%%%%%%%%%%%%%%%%%%%%%%%%%%%%%%%%%%%%%%%%%%%%%%%%%%%%%%%%%
\subsection{Estimation Error}
\label{subsec:financial-estimation-error}
%%%%%%%%%%%%%%%%%%%%%%%%%%%%%%%%%%%%%%%%%%%%%%%%%%%%%%%%%%%%

Portfolio and pricing models can be highly sensitive to estimated parameters.

In mean-variance optimization, small errors in

\[
\boldsymbol{\mu}
\]

may generate large changes in optimal portfolio weights.

Thus parameter uncertainty is often as important as optimization precision.

%%%%%%%%%%%%%%%%%%%%%%%%%%%%%%%%%%%%%%%%%%%%%%%%%%%%%%%%%%%%
\subsection{Model Risk}
\label{subsec:financial-model-risk}
%%%%%%%%%%%%%%%%%%%%%%%%%%%%%%%%%%%%%%%%%%%%%%%%%%%%%%%%%%%%

Model risk arises when the mathematical model does not adequately represent
the financial system.

Possible causes include:

\[
\boxed{
\text{incorrect assumptions},
}
\]

\[
\boxed{
\text{poor calibration},
}
\]

\[
\boxed{
\text{structural breaks},
}
\]

and

\[
\boxed{
\text{unmodeled market behavior}.
}
\]

%%%%%%%%%%%%%%%%%%%%%%%%%%%%%%%%%%%%%%%%%%%%%%%%%%%%%%%%%%%%
\subsection{Stress Testing}
\label{subsec:financial-stress-testing}
%%%%%%%%%%%%%%%%%%%%%%%%%%%%%%%%%%%%%%%%%%%%%%%%%%%%%%%%%%%%

Stress testing evaluates portfolio or institutional performance under severe
scenarios.

Examples include:

\[
\boxed{
\text{large equity decline},
}
\]

\[
\boxed{
\text{interest-rate shock},
}
\]

\[
\boxed{
\text{volatility spike},
}
\]

and

\[
\boxed{
\text{credit deterioration}.
}
\]

Stress tests complement probability-based risk measures.

%%%%%%%%%%%%%%%%%%%%%%%%%%%%%%%%%%%%%%%%%%%%%%%%%%%%%%%%%%%%
\subsection{Scenario Analysis}
\label{subsec:financial-scenario-analysis}
%%%%%%%%%%%%%%%%%%%%%%%%%%%%%%%%%%%%%%%%%%%%%%%%%%%%%%%%%%%%

Scenario analysis evaluates outcomes under specified combinations of future
conditions.

For scenario \(s\),

\[
\boxed{
L_s
=
L(\mathbf{z}_s),
}
\]

where \(\mathbf{z}_s\) contains relevant financial variables.

Scenarios may be:

\[
\boxed{
\text{historical},
}
\]

\[
\boxed{
\text{hypothetical},
}
\]

or

\[
\boxed{
\text{statistically simulated}.
}
\]

%%%%%%%%%%%%%%%%%%%%%%%%%%%%%%%%%%%%%%%%%%%%%%%%%%%%%%%%%%%%
\subsection{Financial Networks}
\label{subsec:mathematical-finance-networks}
%%%%%%%%%%%%%%%%%%%%%%%%%%%%%%%%%%%%%%%%%%%%%%%%%%%%%%%%%%%%

Financial institutions can be represented as nodes.

Edges may represent:

\[
\boxed{
\text{loans},
}
\]

\[
\boxed{
\text{derivative exposures},
}
\]

\[
\boxed{
\text{payment obligations},
}
\]

or

\[
\boxed{
\text{shared asset exposure}.
}
\]

Network structure can influence systemic risk.

%%%%%%%%%%%%%%%%%%%%%%%%%%%%%%%%%%%%%%%%%%%%%%%%%%%%%%%%%%%%
\subsection{Systemic Risk}
\label{subsec:financial-systemic-risk}
%%%%%%%%%%%%%%%%%%%%%%%%%%%%%%%%%%%%%%%%%%%%%%%%%%%%%%%%%%%%

Systemic risk is the possibility that distress spreads across the financial
system.

Mechanisms include:

\[
\boxed{
\text{direct counterparty contagion},
}
\]

\[
\boxed{
\text{liquidity spirals},
}
\]

\[
\boxed{
\text{fire sales},
}
\]

and

\[
\boxed{
\text{common exposures}.
}
\]

Thus systemic risk is not simply the sum of individual institution risks.

%%%%%%%%%%%%%%%%%%%%%%%%%%%%%%%%%%%%%%%%%%%%%%%%%%%%%%%%%%%%
\subsection{High-Frequency Financial Data}
\label{subsec:high-frequency-finance}
%%%%%%%%%%%%%%%%%%%%%%%%%%%%%%%%%%%%%%%%%%%%%%%%%%%%%%%%%%%%

At very short time scales, market microstructure becomes important.

Observed prices are affected by:

\[
\boxed{
\text{bid-ask bounce},
}
\]

\[
\boxed{
\text{discrete price ticks},
}
\]

\[
\boxed{
\text{irregular trading times},
}
\]

and

\[
\boxed{
\text{order-book dynamics}.
}
\]

Continuous diffusion models may be inadequate at these scales.

%%%%%%%%%%%%%%%%%%%%%%%%%%%%%%%%%%%%%%%%%%%%%%%%%%%%%%%%%%%%
\subsection{No-Arbitrage and Reality}
\label{subsec:no-arbitrage-and-reality}
%%%%%%%%%%%%%%%%%%%%%%%%%%%%%%%%%%%%%%%%%%%%%%%%%%%%%%%%%%%%

No-arbitrage is a modeling principle.

Real markets contain:

\[
\boxed{
\text{transaction costs},
}
\]

\[
\boxed{
\text{funding constraints},
}
\]

\[
\boxed{
\text{short-sale restrictions},
}
\]

\[
\boxed{
\text{latency},
}
\]

and

\[
\boxed{
\text{liquidity limitations}.
}
\]

Therefore theoretical arbitrage arguments must be interpreted relative to
the market assumptions.

%%%%%%%%%%%%%%%%%%%%%%%%%%%%%%%%%%%%%%%%%%%%%%%%%%%%%%%%%%%%
\subsection{Mathematical Finance Workflow}
\label{subsec:mathematical-finance-workflow}
%%%%%%%%%%%%%%%%%%%%%%%%%%%%%%%%%%%%%%%%%%%%%%%%%%%%%%%%%%%%

A useful modeling workflow is:

\begin{enumerate}
    \item define the financial question;
    \item identify the assets and contracts;
    \item specify the time horizon;
    \item define the relevant cash flows;
    \item select return and compounding conventions;
    \item identify risk factors;
    \item determine whether a static or dynamic model is needed;
    \item estimate or calibrate model parameters;
    \item specify probability assumptions;
    \item identify tradable hedging instruments;
    \item check for no-arbitrage relationships;
    \item construct the pricing or optimization model;
    \item determine whether markets are modeled as complete or incomplete;
    \item incorporate transaction costs when relevant;
    \item incorporate constraints;
    \item select risk measures;
    \item solve analytically when possible;
    \item otherwise use trees, PDE methods, optimization, or simulation;
    \item test numerical convergence;
    \item perform sensitivity analysis;
    \item evaluate parameter uncertainty;
    \item evaluate model risk;
    \item perform stress tests;
    \item compare model outputs with market data;
    \item clearly state assumptions and limitations.
\end{enumerate}

%%%%%%%%%%%%%%%%%%%%%%%%%%%%%%%%%%%%%%%%%%%%%%%%%%%%%%%%%%%%
\subsection{Common Errors}
\label{subsec:mathematical-finance-common-errors}
%%%%%%%%%%%%%%%%%%%%%%%%%%%%%%%%%%%%%%%%%%%%%%%%%%%%%%%%%%%%

\subsubsection{Confusing Simple and Log Returns}

Simple return is

\[
\frac{S_{t+1}-S_t}{S_t}.
\]

Log return is

\[
\ln
\left(
\frac{S_{t+1}}{S_t}
\right).
\]

They are related but not identical.

%%%%%%%%%%%%%%%%%%%%%%%%%%%%%%%%%%%%%%%%%%%%%%%%%%%%%%%%%%%%
\subsubsection{Adding Simple Returns Across Time}

Log returns add exactly across consecutive periods.

Simple returns compound multiplicatively:

\[
\boxed{
1+R_{\mathrm{total}}
=
\prod_t(1+R_t).
}
\]

%%%%%%%%%%%%%%%%%%%%%%%%%%%%%%%%%%%%%%%%%%%%%%%%%%%%%%%%%%%%
\subsubsection{Mixing Compounding Conventions}

An annual effective rate and a continuously compounded rate are not the same
numerical quantity.

Compounding conventions must be kept consistent.

%%%%%%%%%%%%%%%%%%%%%%%%%%%%%%%%%%%%%%%%%%%%%%%%%%%%%%%%%%%%
\subsubsection{Confusing Price and Return}

A high-priced stock does not necessarily have a higher expected return.

Price level and percentage return are different concepts.

%%%%%%%%%%%%%%%%%%%%%%%%%%%%%%%%%%%%%%%%%%%%%%%%%%%%%%%%%%%%
\subsubsection{Ignoring Covariance}

Portfolio risk cannot generally be computed from individual asset
volatilities alone.

Covariances matter:

\[
\boxed{
\sigma_p^2
=
\mathbf{w}^T
\Sigma
\mathbf{w}.
}
\]

%%%%%%%%%%%%%%%%%%%%%%%%%%%%%%%%%%%%%%%%%%%%%%%%%%%%%%%%%%%%
\subsubsection{Assuming More Assets Always Mean Better Diversification}

Diversification depends on dependence structure.

Highly correlated assets may provide little risk reduction.

%%%%%%%%%%%%%%%%%%%%%%%%%%%%%%%%%%%%%%%%%%%%%%%%%%%%%%%%%%%%
\subsubsection{Treating Expected Returns as Known Constants}

Expected return estimates can be extremely uncertain.

Optimization based on noisy estimates may produce unstable portfolios.

%%%%%%%%%%%%%%%%%%%%%%%%%%%%%%%%%%%%%%%%%%%%%%%%%%%%%%%%%%%%
\subsubsection{Ignoring Short-Sale Constraints}

An unconstrained optimizer may return negative weights.

These correspond to short positions and may be prohibited or expensive.

%%%%%%%%%%%%%%%%%%%%%%%%%%%%%%%%%%%%%%%%%%%%%%%%%%%%%%%%%%%%
\subsubsection{Treating Volatility as the Same as All Risk}

Volatility measures dispersion.

It does not by itself capture:

\[
\boxed{
\text{liquidity risk},
}
\]

\[
\boxed{
\text{credit risk},
}
\]

\[
\boxed{
\text{tail risk},
}
\]

or

\[
\boxed{
\text{model risk}.
}
\]

%%%%%%%%%%%%%%%%%%%%%%%%%%%%%%%%%%%%%%%%%%%%%%%%%%%%%%%%%%%%
\subsubsection{Confusing Correlation With Causation}

Financial returns may be correlated because of shared risk factors.

Correlation alone does not establish a causal economic mechanism.

%%%%%%%%%%%%%%%%%%%%%%%%%%%%%%%%%%%%%%%%%%%%%%%%%%%%%%%%%%%%
\subsubsection{Assuming CAPM Is an Exact Law}

CAPM depends on idealized assumptions.

It is a theoretical benchmark rather than a guaranteed empirical equation.

%%%%%%%%%%%%%%%%%%%%%%%%%%%%%%%%%%%%%%%%%%%%%%%%%%%%%%%%%%%%
\subsubsection{Confusing Arbitrage With High Expected Return}

An arbitrage is not merely a profitable investment.

It is a strategy satisfying strict no-loss and no-net-investment conditions
within the model.

%%%%%%%%%%%%%%%%%%%%%%%%%%%%%%%%%%%%%%%%%%%%%%%%%%%%%%%%%%%%
\subsubsection{Using Real-World Probabilities in Risk-Neutral Pricing}

Risk-neutral probabilities are pricing probabilities.

They need not equal estimated real-world probabilities.

%%%%%%%%%%%%%%%%%%%%%%%%%%%%%%%%%%%%%%%%%%%%%%%%%%%%%%%%%%%%
\subsubsection{Interpreting Risk-Neutrality as Investor Risk Preference}

Risk-neutral valuation does not require every investor to be risk-neutral.

It follows from no-arbitrage and replication under suitable assumptions.

%%%%%%%%%%%%%%%%%%%%%%%%%%%%%%%%%%%%%%%%%%%%%%%%%%%%%%%%%%%%
\subsubsection{Using Black--Scholes Without Checking Assumptions}

The basic Black--Scholes model assumes constant volatility, continuous
trading, and other idealizations.

Observed markets violate these assumptions to varying degrees.

%%%%%%%%%%%%%%%%%%%%%%%%%%%%%%%%%%%%%%%%%%%%%%%%%%%%%%%%%%%%
\subsubsection{Confusing Historical and Implied Volatility}

Historical volatility is estimated from past returns.

Implied volatility is inferred from current option prices through a pricing
model.

%%%%%%%%%%%%%%%%%%%%%%%%%%%%%%%%%%%%%%%%%%%%%%%%%%%%%%%%%%%%
\subsubsection{Assuming Implied Volatility Is Constant Across Strikes}

Real option markets typically exhibit smiles or skews.

A single constant volatility generally cannot fit all traded options.

%%%%%%%%%%%%%%%%%%%%%%%%%%%%%%%%%%%%%%%%%%%%%%%%%%%%%%%%%%%%
\subsubsection{Treating Delta Hedging as Perfect in Real Markets}

Continuous exact rebalancing is impossible.

Real hedging contains:

\[
\boxed{
\text{discrete rebalancing error},
}
\]

\[
\boxed{
\text{transaction costs},
}
\]

and

\[
\boxed{
\text{model error}.
}
\]

%%%%%%%%%%%%%%%%%%%%%%%%%%%%%%%%%%%%%%%%%%%%%%%%%%%%%%%%%%%%
\subsubsection{Ignoring Early Exercise in American Options}

European and American contracts have different exercise rights.

American-option valuation requires comparison between continuation and
exercise values.

%%%%%%%%%%%%%%%%%%%%%%%%%%%%%%%%%%%%%%%%%%%%%%%%%%%%%%%%%%%%
\subsubsection{Confusing VaR With Maximum Possible Loss}

VaR is a quantile.

Losses larger than VaR remain possible.

%%%%%%%%%%%%%%%%%%%%%%%%%%%%%%%%%%%%%%%%%%%%%%%%%%%%%%%%%%%%
\subsubsection{Ignoring Tail Severity Beyond VaR}

Two portfolios can have the same VaR but very different extreme-loss
distributions.

CVaR or stress testing can provide additional information.

%%%%%%%%%%%%%%%%%%%%%%%%%%%%%%%%%%%%%%%%%%%%%%%%%%%%%%%%%%%%
\subsubsection{Ignoring Interest-Rate Term Structure}

A single constant rate may be inappropriate when cash flows occur at
different maturities.

Discounting should reflect the relevant term structure.

%%%%%%%%%%%%%%%%%%%%%%%%%%%%%%%%%%%%%%%%%%%%%%%%%%%%%%%%%%%%
\subsubsection{Ignoring Credit Risk in Discounting}

A risky bond is not generally priced by simply discounting all promised cash
flows at a risk-free rate.

Default and recovery must be considered.

%%%%%%%%%%%%%%%%%%%%%%%%%%%%%%%%%%%%%%%%%%%%%%%%%%%%%%%%%%%%
\subsubsection{Ignoring Transaction Costs in Frequent Trading Strategies}

A strategy profitable before costs can become unprofitable after:

\[
\boxed{
\text{spreads},
}
\]

\[
\boxed{
\text{commissions},
}
\]

and

\[
\boxed{
\text{market impact}.
}
\]

%%%%%%%%%%%%%%%%%%%%%%%%%%%%%%%%%%%%%%%%%%%%%%%%%%%%%%%%%%%%
\subsubsection{Overfitting Financial Models}

A model may fit historical data extremely well while performing poorly
out-of-sample.

Financial systems change over time.

%%%%%%%%%%%%%%%%%%%%%%%%%%%%%%%%%%%%%%%%%%%%%%%%%%%%%%%%%%%%
\subsubsection{Ignoring Model Risk}

Precise numerical outputs do not imply the model itself is correct.

Sensitivity to alternative models should be examined.

%%%%%%%%%%%%%%%%%%%%%%%%%%%%%%%%%%%%%%%%%%%%%%%%%%%%%%%%%%%%
\subsection{Applications}
\label{subsec:mathematical-finance-applications}
%%%%%%%%%%%%%%%%%%%%%%%%%%%%%%%%%%%%%%%%%%%%%%%%%%%%%%%%%%%%

\begin{applicationbox}

\textbf{Portfolio Management.}
Optimization balances expected return, covariance, constraints, and risk
preferences.

\medskip

\textbf{Derivative Pricing.}
No-arbitrage, stochastic calculus, PDEs, trees, and simulation determine
model prices of options and other contingent claims.

\medskip

\textbf{Risk Management.}
VaR, CVaR, stress testing, scenario analysis, and sensitivity measures
quantify financial exposures.

\medskip

\textbf{Fixed Income.}
Discount factors, yield curves, duration, convexity, and interest-rate models
describe bonds and rate-sensitive securities.

\medskip

\textbf{Credit Risk.}
Probability models estimate default, loss, recovery, and counterparty
exposure.

\medskip

\textbf{Asset Allocation.}
Dynamic and stochastic optimization determine investment policies over time.

\medskip

\textbf{Algorithmic Trading.}
Statistical models, stochastic control, optimization, and market
microstructure influence automated trading strategies.

\medskip

\textbf{Systemic Risk.}
Network models describe contagion and interconnected exposure across
financial institutions.

\end{applicationbox}

%%%%%%%%%%%%%%%%%%%%%%%%%%%%%%%%%%%%%%%%%%%%%%%%%%%%%%%%%%%%
\subsection{Exercises}
\label{subsec:mathematical-finance-exercises}
%%%%%%%%%%%%%%%%%%%%%%%%%%%%%%%%%%%%%%%%%%%%%%%%%%%%%%%%%%%%

\begin{exercise}[1]
Define a financial asset.
\end{exercise}

\begin{exercise}[1]
Define simple return.
\end{exercise}

\begin{exercise}[1]
Define logarithmic return.
\end{exercise}

\begin{exercise}[1]
Define present value.
\end{exercise}

\begin{exercise}[1]
Define a zero-coupon bond.
\end{exercise}

\begin{exercise}[1]
Define diversification.
\end{exercise}

\begin{exercise}[1]
Define arbitrage.
\end{exercise}

\begin{exercise}[1]
Define a European call.
\end{exercise}

\begin{exercise}[1]
Define a European put.
\end{exercise}

\begin{exercise}[1]
Define volatility.
\end{exercise}

\begin{exercise}[2]
Compute the simple and log returns when price rises from \(80\) to \(86\).
\end{exercise}

\begin{exercise}[2]
Show that log returns add over consecutive periods.
\end{exercise}

\begin{exercise}[2]
Compute the future value of \(1000\) invested for \(5\) years at a
continuously compounded rate of \(4\%\).
\end{exercise}

\begin{exercise}[2]
Compute the present value of \(5000\) received in \(4\) years at
continuously compounded rate \(6\%\).
\end{exercise}

\begin{exercise}[2]
Price a zero-coupon bond paying \(100\) in \(3\) years at constant
continuously compounded rate \(5\%\).
\end{exercise}

\begin{exercise}[3]
Price a coupon bond by discounting each cash flow.
\end{exercise}

\begin{exercise}[3]
Explain why bond price decreases when yield increases.
\end{exercise}

\begin{exercise}[2]
Explain duration conceptually.
\end{exercise}

\begin{exercise}[2]
Explain convexity conceptually.
\end{exercise}

\begin{exercise}[2]
Compute expected return of a two-asset portfolio.
\end{exercise}

\begin{exercise}[2]
Compute variance of a two-asset portfolio.
\end{exercise}

\begin{exercise}[3]
Analyze how portfolio variance changes as correlation varies from \(-1\) to
\(1\).
\end{exercise}

\begin{exercise}[3]
Show mathematically why imperfect correlation can produce diversification.
\end{exercise}

\begin{exercise}[2]
Write the matrix formula for portfolio variance.
\end{exercise}

\begin{exercise}[3]
Derive the global minimum-variance portfolio.
\end{exercise}

\begin{exercise}[3]
Formulate a minimum-variance portfolio with no short selling.
\end{exercise}

\begin{exercise}[2]
Define the efficient frontier.
\end{exercise}

\begin{exercise}[2]
Compute the Sharpe ratio of a portfolio.
\end{exercise}

\begin{exercise}[2]
Compute market beta from covariance and variance.
\end{exercise}

\begin{exercise}[3]
Use CAPM to compute a model expected return.
\end{exercise}

\begin{exercise}[3]
Explain why CAPM should not be interpreted as an exact empirical law.
\end{exercise}

\begin{exercise}[2]
Write the payoff of a long forward contract.
\end{exercise}

\begin{exercise}[2]
Compute a forward price using

\[
F_0=S_0e^{rT}.
\]
\end{exercise}

\begin{exercise}[3]
Construct a cash-and-carry argument for the forward-price formula.
\end{exercise}

\begin{exercise}[2]
Graph the payoff of a European call.
\end{exercise}

\begin{exercise}[2]
Graph the payoff of a European put.
\end{exercise}

\begin{exercise}[2]
State put-call parity.
\end{exercise}

\begin{exercise}[3]
Derive put-call parity using two portfolios.
\end{exercise}

\begin{exercise}[3]
Find a missing option price using put-call parity.
\end{exercise}

\begin{exercise}[2]
Construct a one-period binomial tree.
\end{exercise}

\begin{exercise}[3]
Compute the binomial hedge ratio.
\end{exercise}

\begin{exercise}[3]
Compute the risk-neutral probability.
\end{exercise}

\begin{exercise}[3]
Price a European call in a one-period binomial model.
\end{exercise}

\begin{exercise}[3]
Price a European put in a multi-period binomial tree.
\end{exercise}

\begin{exercise}[3]
Explain backward induction.
\end{exercise}

\begin{exercise}[3]
Explain early exercise in an American option.
\end{exercise}

\begin{exercise}[2]
Write the geometric Brownian motion model.
\end{exercise}

\begin{exercise}[3]
Apply It\^o's formula to derive the SDE for \(\ln S_t\).
\end{exercise}

\begin{exercise}[2]
Explain historical volatility.
\end{exercise}

\begin{exercise}[2]
Explain implied volatility.
\end{exercise}

\begin{exercise}[3]
Explain why implied volatility can vary with strike.
\end{exercise}

\begin{exercise}[2]
Explain risk-neutral valuation.
\end{exercise}

\begin{exercise}[3]
Explain why risk-neutral probabilities need not equal real-world
probabilities.
\end{exercise}

\begin{exercise}[3]
Explain why the drift \(\mu\) disappears from the Black--Scholes PDE.
\end{exercise}

\begin{exercise}[2]
State the Black--Scholes PDE.
\end{exercise}

\begin{exercise}[2]
State the Black--Scholes call formula.
\end{exercise}

\begin{exercise}[2]
Define delta.
\end{exercise}

\begin{exercise}[2]
Define gamma.
\end{exercise}

\begin{exercise}[2]
Define theta.
\end{exercise}

\begin{exercise}[2]
Define vega.
\end{exercise}

\begin{exercise}[2]
Define rho.
\end{exercise}

\begin{exercise}[3]
Explain delta hedging.
\end{exercise}

\begin{exercise}[3]
Explain why delta hedging requires rebalancing.
\end{exercise}

\begin{exercise}[3]
Explain the role of gamma in hedging error.
\end{exercise}

\begin{exercise}[2]
Write a Monte Carlo estimator for a discounted derivative payoff.
\end{exercise}

\begin{exercise}[3]
Explain why Monte Carlo convergence is proportional to \(N^{-1/2}\).
\end{exercise}

\begin{exercise}[3]
Explain a control-variate method for option pricing.
\end{exercise}

\begin{exercise}[3]
Explain antithetic variates.
\end{exercise}

\begin{exercise}[2]
Explain finite-difference option pricing conceptually.
\end{exercise}

\begin{exercise}[3]
State the terminal condition for a European call in the Black--Scholes PDE.
\end{exercise}

\begin{exercise}[2]
Define stochastic volatility conceptually.
\end{exercise}

\begin{exercise}[2]
Define a jump-diffusion model conceptually.
\end{exercise}

\begin{exercise}[3]
Explain why jumps may improve modeling of extreme short-term returns.
\end{exercise}

\begin{exercise}[2]
Define a yield curve.
\end{exercise}

\begin{exercise}[2]
Compute a continuously compounded spot rate from a discount factor.
\end{exercise}

\begin{exercise}[3]
Derive the instantaneous forward-rate formula.
\end{exercise}

\begin{exercise}[3]
Explain mean reversion in interest rates.
\end{exercise}

\begin{exercise}[2]
Define credit risk.
\end{exercise}

\begin{exercise}[2]
Compute expected credit loss from PD, LGD, and EAD.
\end{exercise}

\begin{exercise}[2]
Define Value at Risk.
\end{exercise}

\begin{exercise}[3]
Explain why VaR is not a maximum-loss measure.
\end{exercise}

\begin{exercise}[2]
Define Conditional Value at Risk.
\end{exercise}

\begin{exercise}[3]
Compare VaR and CVaR.
\end{exercise}

\begin{exercise}[3]
Formulate a portfolio optimization problem with a CVaR constraint.
\end{exercise}

\begin{exercise}[2]
Explain transaction costs.
\end{exercise}

\begin{exercise}[2]
Define the bid-ask spread.
\end{exercise}

\begin{exercise}[2]
Explain liquidity risk.
\end{exercise}

\begin{exercise}[3]
Explain market impact.
\end{exercise}

\begin{exercise}[3]
Formulate an optimal trade problem with quadratic market impact.
\end{exercise}

\begin{exercise}[2]
Explain calibration.
\end{exercise}

\begin{exercise}[3]
Formulate a least-squares calibration problem.
\end{exercise}

\begin{exercise}[3]
Explain parameter-estimation risk.
\end{exercise}

\begin{exercise}[2]
Define model risk.
\end{exercise}

\begin{exercise}[2]
Define stress testing.
\end{exercise}

\begin{exercise}[3]
Construct a multi-factor financial stress scenario.
\end{exercise}

\begin{exercise}[3]
Explain systemic risk using a network of institutions.
\end{exercise}

\begin{exercise}[4]
Derive the mean-variance efficient frontier analytically.
\end{exercise}

\begin{exercise}[4]
Study portfolio optimization under short-sale constraints.
\end{exercise}

\begin{exercise}[4]
Study sensitivity of optimal weights to errors in expected returns.
\end{exercise}

\begin{exercise}[4]
Derive the one-period binomial pricing formula by replication.
\end{exercise}

\begin{exercise}[4]
Extend binomial pricing to multiple periods.
\end{exercise}

\begin{exercise}[4]
Price an American put by backward induction.
\end{exercise}

\begin{exercise}[4]
Derive the Black--Scholes PDE using a delta-hedged portfolio.
\end{exercise}

\begin{exercise}[4]
Derive the Black--Scholes European call formula.
\end{exercise}

\begin{exercise}[4]
Verify put-call parity using Black--Scholes formulas.
\end{exercise}

\begin{exercise}[4]
Compute and graph option Greeks as functions of \(S\).
\end{exercise}

\begin{exercise}[4]
Simulate geometric Brownian motion.
\end{exercise}

\begin{exercise}[4]
Price a European option by Monte Carlo simulation.
\end{exercise}

\begin{exercise}[4]
Compare Monte Carlo and Black--Scholes prices.
\end{exercise}

\begin{exercise}[4]
Implement variance reduction for Monte Carlo option pricing.
\end{exercise}

\begin{exercise}[4]
Solve the Black--Scholes PDE numerically by finite differences.
\end{exercise}

\begin{exercise}[4]
Calibrate implied volatility from option prices.
\end{exercise}

\begin{exercise}[4]
Construct an implied-volatility smile.
\end{exercise}

\begin{exercise}[4]
Study a stochastic-volatility model numerically.
\end{exercise}

\begin{exercise}[4]
Study a jump-diffusion option-pricing model.
\end{exercise}

\begin{exercise}[4]
Calibrate a simple yield curve from zero-coupon bond prices.
\end{exercise}

\begin{exercise}[4]
Simulate a mean-reverting short-rate process.
\end{exercise}

\begin{exercise}[4]
Compute portfolio VaR by historical simulation.
\end{exercise}

\begin{exercise}[4]
Compute portfolio CVaR by Monte Carlo simulation.
\end{exercise}

\begin{exercise}[4]
Optimize a portfolio using CVaR.
\end{exercise}

\begin{exercise}[4]
Study dynamic portfolio allocation as a stochastic-control problem.
\end{exercise}

\begin{exercise}[4]
Analyze the effect of transaction costs on delta hedging.
\end{exercise}

\begin{exercise}[4]
Construct a financial contagion model on a network.
\end{exercise}

%%%%%%%%%%%%%%%%%%%%%%%%%%%%%%%%%%%%%%%%%%%%%%%%%%%%%%%%%%%%
\subsection{Conceptual Summary}
\label{subsec:mathematical-finance-summary}
%%%%%%%%%%%%%%%%%%%%%%%%%%%%%%%%%%%%%%%%%%%%%%%%%%%%%%%%%%%%

Mathematical finance applies probability, optimization, stochastic calculus,
and numerical mathematics to financial assets and decisions.

Returns may be measured by

\[
\boxed{
R
=
\frac{S_1-S_0}{S_0}
}
\]

or

\[
\boxed{
r
=
\ln\frac{S_1}{S_0}.
}
\]

Discounting converts future cash flows to present value:

\[
\boxed{
PV
=
Ce^{-rT}.
}
\]

Portfolio return is

\[
\boxed{
R_p
=
\mathbf{w}^T\mathbf{R},
}
\]

while portfolio variance is

\[
\boxed{
\sigma_p^2
=
\mathbf{w}^T
\Sigma
\mathbf{w}.
}
\]

Diversification arises from imperfect dependence among asset returns.

Derivative pricing is built on

\[
\boxed{
\text{no arbitrage}
}
\]

and

\[
\boxed{
\text{replication}.
}
\]

A fundamental stochastic asset model is

\[
\boxed{
dS_t
=
\mu S_tdt
+
\sigma S_tdW_t.
}
\]

Under risk-neutral valuation,

\[
\boxed{
V_t
=
\mathbb{E}^{Q}
\left[
e^{-r(T-t)}
H_T
\mid
\mathcal{F}_t
\right].
}
\]

For European derivatives in the Black--Scholes framework,

\[
\boxed{
V_t
+
\frac12\sigma^2S^2V_{SS}
+
rSV_S
-
rV
=
0.
}
\]

Risk management uses tools such as

\[
\boxed{
\operatorname{VaR},
}
\]

\[
\boxed{
\operatorname{CVaR},
}
\]

\[
\boxed{
\text{stress tests},
}
\]

and

\[
\boxed{
\text{scenario analysis}.
}
\]

Mathematical finance therefore combines

\[
\boxed{
\text{probability}
+
\text{optimization}
+
\text{stochastic calculus}
+
\text{PDEs}
+
\text{statistics}
+
\text{computation}
}
\]

to study pricing, portfolios, hedging, risk, and financial decision making.

%%%%%%%%%%%%%%%%%%%%%%%%%%%%%%%%%%%%%%%%%%%%%%%%%%%%%%%%%%%%
\subsection{Section Connections}
\label{subsec:mathematical-finance-connections}
%%%%%%%%%%%%%%%%%%%%%%%%%%%%%%%%%%%%%%%%%%%%%%%%%%%%%%%%%%%%

\chapterconnection{
Mathematical Finance applies probability, stochastic processes, stochastic
calculus, optimization, partial differential equations, and numerical
methods to asset prices, portfolios, derivatives, and financial risk. The
portfolio models connect directly with Optimization and Operations Research,
continuous-time asset models connect with Stochastic Differential Equations,
option pricing connects with PDEs and Monte Carlo methods, dynamic portfolio
choice connects with Control Theory, and systemic financial risk connects
with Network Science. The next section, Mathematical Economics, develops
mathematical models of consumer and producer behavior, markets, equilibrium,
utility, production, optimization, game theory, growth, and dynamic economic
systems.
}

%%%%%%%%%%%%%%%%%%%%%%%%%%%%%%%%%%%%%%%%%%%%%%%%%%%%%%%%%%%%
% SECTION SOURCE TRACKING
%%%%%%%%%%%%%%%%%%%%%%%%%%%%%%%%%%%%%%%%%%%%%%%%%%%%%%%%%%%%

%------------------------------------------------------------
% 13.8 Mathematical Finance
%------------------------------------------------------------
%
% Source basis:
%   - Standard mathematical finance.
%   - Standard portfolio theory.
%   - Standard no-arbitrage and derivative-pricing theory.
%   - Standard stochastic asset-price models and financial risk methods.
%
% Used for:
%   - financial assets;
%   - simple and logarithmic returns;
%   - compounding and discounting;
%   - present and future value;
%   - discount factors;
%   - zero-coupon and coupon bonds;
%   - yield to maturity;
%   - duration and convexity;
%   - portfolios;
%   - expected return;
%   - covariance and portfolio variance;
%   - diversification;
%   - mean-variance optimization;
%   - minimum-variance portfolios;
%   - efficient frontiers;
%   - risk-free assets;
%   - Sharpe ratios;
%   - factor models;
%   - beta;
%   - CAPM;
%   - arbitrage;
%   - law of one price;
%   - replication;
%   - complete and incomplete markets;
%   - forwards;
%   - forward prices;
%   - calls and puts;
%   - option payoffs;
%   - put-call parity;
%   - binomial pricing;
%   - risk-neutral probabilities;
%   - American options;
%   - geometric Brownian motion;
%   - log-price dynamics;
%   - volatility;
%   - historical and implied volatility;
%   - risk-neutral valuation;
%   - Black--Scholes model;
%   - Black--Scholes PDE;
%   - European call and put formulas;
%   - option Greeks;
%   - delta hedging;
%   - dynamic hedging;
%   - volatility smiles and surfaces;
%   - Monte Carlo pricing;
%   - variance reduction;
%   - finite-difference pricing;
%   - stochastic volatility;
%   - jump diffusion;
%   - term structures;
%   - spot and forward rates;
%   - short-rate models;
%   - credit risk;
%   - expected credit loss;
%   - market risk;
%   - Value at Risk;
%   - Conditional Value at Risk;
%   - risk-constrained optimization;
%   - dynamic portfolio choice;
%   - transaction costs;
%   - bid-ask spreads;
%   - liquidity risk;
%   - market impact;
%   - calibration;
%   - parameter estimation;
%   - estimation error;
%   - model risk;
%   - stress testing;
%   - scenario analysis;
%   - financial networks;
%   - systemic risk;
%   - high-frequency data;
%   - applications and exercises.
%
% Notes:
%   - Mathematical Economics is developed next in Section 13.9.
%
%%%%%%%%%%%%%%%%%%%%%%%%%%%%%%%%%%%%%%%%%%%%%%%%%%%%%%%%%%%%